\documentclass[12pt,twoside]{article}


%%%%%%%%%%%%%%%%%%%%%%%%%%%%%%%%%%%%%%%%%%
 \usepackage{float}
\usepackage{graphicx}
\usepackage{epstopdf}
\usepackage{graphicx}
\usepackage{epic}
\usepackage{multirow}
\usepackage{pst-poly}  %»­Í¼
\usepackage{pst-plot}  %×ø±ê
\usepackage{pst-poly}  %¶à±ßÐÎ
\usepackage{tikz}
\usepackage{threeparttable}
\usepackage{xcolor}
\usetikzlibrary{arrows,shapes,chains}


\renewcommand{\paragraph}{\roman{paragraph}}
\usepackage[a4paper]{geometry}
\setlength{\textwidth}{6.3in}
\setlength{\textheight}{8.8in}
\setlength{\topmargin}{0pt}
\setlength{\headsep}{25pt}
\setlength{\headheight}{0pt}
\setlength{\oddsidemargin}{0pt}
\setlength{\evensidemargin}{0pt}

%%%%%%%%%%%%%%%%%%%%%%%%%%%%%%%%%%%%%%%%%%

\makeatletter
\renewcommand\title[1]{\gdef\@title{\reset@font\Large\bfseries #1}}
\renewcommand\section{\@startsection {section}{1}{\z@}%
                                   {-3.5ex \@plus -1ex \@minus -.2ex}%
                                   {2.3ex \@plus.2ex}%
                                   {\normalfont\large\bfseries}}
\renewcommand\subsection{\@startsection{subsection}{2}{\z@}%
                                     {-3ex\@plus -1ex \@minus -.2ex}%
                                     {1.5ex \@plus .2ex}%
                                     {\normalfont\normalsize\bfseries}}
\renewcommand\subsubsection{\@startsection{subsubsection}{3}{\z@}%
                                     {-2.5ex\@plus -1ex \@minus -.2ex}%
                                     {1.5ex \@plus .2ex}%
                                     {\normalfont\normalsize\bfseries}}

\def\@runningauthor{}\newcommand{\runningauthor}[1]{\def\runningauthor{#1}}
\def\@runningtitle{}\newcommand{\runningtitle}[1]{\def\runningtitle{#1}}

\renewcommand{\ps@plain}{%
\renewcommand{\@evenhead}{\footnotesize\scshape \hfill\runningauthor\hfill}
\renewcommand{\@oddhead}{\footnotesize\scshape \hfill\runningtitle\hfill}}
\newcommand{\Z}{\mathbb{Z}}
\newcommand{\F}{\mathbb{F}}
\newcommand{\x}{\mathbf{x}}
\newcommand{\A}{\mathbf{a}}
\newcommand {\R}{\ms{R}}
\newcommand {\B}{\ms{B}}
\newcommand {\per}{{\em {per}}}
\newcommand {\C}{{\mathcal{C}}}
\newcommand {\PP}{{\mathcal{P}}}
\newcommand {\BB}{{\mathcal{B}}}
\newcommand {\ccc}{{\mathbf{c}}}
\pagestyle{plain}


\g@addto@macro\bfseries{\boldmath}

\makeatother

%%%%%%%%%%%%%%%%%%%%%%%%%%%%%%%%%%%%%%%%%%%%%%%%%%%%%%%%%%%%%%%%%%%%%%%%

% Please remove all other commands that change parameterss such as
% margins or pagesizes.

% we recommend these ams packages
\usepackage{amsthm,amsmath,amssymb}
\usepackage{cite}
% we recommend the graphicx package for importing figures
\usepackage{graphicx}

% hyperlinks
\usepackage[colorlinks=true,citecolor=black,linkcolor=black,urlcolor=blue]{hyperref}

% declare theorem-like environments
\theoremstyle{plain}
\newtheorem{theorem}{Theorem}[section]
\newtheorem{notation}{Notation}
\newtheorem{lem}[theorem]{Lemma}
\newtheorem{cor}[theorem]{Corollary}
\newtheorem{prop}[theorem]{Proposition}
\newtheorem{fact}[theorem]{Fact}
\newtheorem{observation}[theorem]{Observation}
\newtheorem{claim}[theorem]{Claim}

\theoremstyle{definition}
\newtheorem{definition}[theorem]{Definition}
\newtheorem{example}[theorem]{Example}
\newtheorem{conjecture}[theorem]{Conjecture}
\newtheorem{open}[theorem]{Open Problem}
\newtheorem{problem}[theorem]{Problem}
\newtheorem{question}[theorem]{Question}

\theoremstyle{remark}
\newtheorem{remark}[theorem]{Remark}
\newtheorem{note}[theorem]{Note}

%%%%%%%%%%%%%%%%%%%%%%%%%%%%%%%%%%%%%%%%%%%%%%%%%%%%%%%

%% Enter a title


%% Enter a running title

%\runningtitle{The $q$-ary antiprimitive BCH codes}

%% Enter authors and affiliations

%\author{Hongwei Zhu and Minjia Shi\thanks{Key Laboratory of Intelligent Computing and Signal Processing, Ministry of Education, School of Mathematical Sciences, Anhui University, Hefei, 230601, China. E-mail: zhwgood66@163.com, smjwcl.good@163.com}
%\and Xiaoqiang Wang\thanks{Hubei Key Laboratory of Applied Mathematics, Faculty of Mathematics and Statistics, Hubei University, Wuhan 430062, China.
%email: waxiqq@163.com}
%\and Tor Helleseth\thanks{Department of Informatics, University of Bergen, Bergen, Norway. E-mail: tor.helleseth@uib.no}\\
%\small School of Mathematical Sciences\\[-0.8ex]
%\small Anhui University, Anhui, Hefei\\[-0.8ex]
%\small\tt smjwcl.good@163.com\\
%\and
%Patrick Sol\'e \\
%\small 2 rue de la Libert\'e, 93 526 Saint-Denis, France\\[-0.8ex]
%\small CNRS/LAGA, University of Paris 8\\[-0.8ex]
%\small\tt sole@enst.fr
%}

%% Enter running author names

\runningauthor{}

\date{}


\begin{document}


  \title{Improved lower and upper bounds for LCD codes\thanks{This research is supported by the National Natural Science Foundation of China (12071001) and 2021 University Graduate Research Project (Y020410077).}}
\author{ Shitao Li\thanks{lishitao0216@163.com}, Minjia Shi\thanks{smjwcl.good@163.com}
\thanks{ Shitao Li and Minjia Shi are with School of Mathematical Sciences, Anhui University, Hefei, China.}}

\date{}
    \maketitle

\begin{abstract}
Linear complementary dual (LCD) codes are linear codes which intersect their dual codes trivially, which have been
of interest and extensively studied due to its wide applications.
In this paper, we give some methods for constructing LCD codes over small fields by modifying some known methods. We show that all odd-like binary LCD codes, ternary LCD codes and quaternary Hermitian LCD codes can be constructed by the modified methods. Using these methods, we construct a lot of optimal binary LCD codes, ternary LCD codes and quaternary Hermitian LCD codes, which improve the known lower bounds on the largest minimum weights. Furthermore, we give two counterexamples to show that the conjecture proposed by Bouyuklieva (Des. Codes Cryptogr. 89(11): 2445-2461, 2021) is invalid.
\end{abstract}
{\bf Keywords:} Binary LCD codes, ternary LCD codes, quaternary Hermitian LCD codes\\
{\bf Mathematics Subject Classification} 94B05 15B05 12E10
\section{Introduction}
LCD codes were first introduced by Massey in 1992 to solve a problem in information theory \cite{LCD-7}.
In 2004, Sendrier \cite{LCD-is-good} showed that LCD codes meet the asymptotic Gilbert-Varshamov bound by using the hull dimension spectra of linear codes.
In 2014, Carlet {\it{et al.}} \cite{lcd-appl} investigated an application of binary LCD codes against Side-Channel Attacks (SCA) and Fault Injection Attack (FIA), and gave several constructions of LCD codes.
Recently, LCD codes were extensively studied \cite{LCD-1,LCD-new,LCD-2,LCD-3,
LCD-5,LCD-9,LCD-huang,ACD-E-codes,LCD-T-matric,LCD-zhu}.
In particular, Carlet {\it{et al.}} \cite{LCD-1} showed that any code
over $\F_q$ is equivalent to some Euclidean LCD code for $q\geq 4$ and any code over $\F_{q^2}$ is
equivalent to some Hermitian LCD code for $q\geq 3$. This motivates us to study
LCD codes in small fields.

It is also a fundamental topic to determine the largest minimum distance of LCD codes for various lengths and dimensions in coding
theory.
In recent years, much work has been done concerning this fundamental topic.
The largest minimum weights among all binary LCD $[n,k]$ codes were partially determined
in \cite{bound-12,HS-BLCD-1-16,2-LCD-30,AH-BLCD-17-24,B-LCD-40,2-LCD-40} for $n\leq 40$. It is worth noting that there are many unknowns in \cite{B-LCD-40} and \cite{2-LCD-40}. The largest minimum weights among all ternary LCD $[n,k]$ codes were determined in \cite{AH-TLCD-1-10,ter-11-19} for $n\leq 20$. The largest minimum weights among all quaternary Hermitian LCD $[n,k]$ codes were partially determined in \cite{H-LCD-20,H-20,H-25} for $n\leq 25$.
More constructions of LCD codes can be seen \cite{H-LCD-IT,H-lCD-k=3,AHS-BLCD,H-lcd-IJQI,Lin-sok,LCD-10}.

Furthermore, Carlet {\it{et al.}} \cite{LCD-1} also proved that any linear code is the direct sum of a self-orthogonal code and an LCD code with respect to the Euclidean inner product. As a result, they presented a new approach of constructing LCD codes by extending linear codes.  Harada \cite{234-lcD} gave two methods for constructing many LCD codes from a given LCD code by modifying some known methods for constructing self-dual codes \cite{b-up-1,b-up-2,b-up-3,b-up-4,b-up-5,b-up-6,b-up-7}.
Using these methods, many new binary LCD codes and quaternary Hermitian LCD were constructed.

In this paper, based on ideas of \cite{LCD-1} and \cite{234-lcD}, we further modify these methods. We show that all odd-like binary LCD codes, ternary LCD codes and quaternary Hermitian LCD codes can be constructed by our construction methods. That is to say, our methods are efficient for constructing LCD codes.
Using theses methods,
we obtain that some binary LCD codes with better parameters comparing with \cite{B-LCD-40} and \cite{2-LCD-40}. In addition, we also obtain some binary LCD codes, which are not equivalent to the codes in \cite{2-LCD-40}.
We extend the tables on ternary LCD codes to length 25. Some ternary LCD codes with better parameters are constructed comparing with \cite{234-lcD}. Finally, we also construct some quaternary Hermitian LCD codes with new parameters comparing with \cite{H-25}. These codes improve the previously known lower bounds on the largest minimum weights.
It is worth mentioning that we give two counterexamples to show that the conjecture proposed by Bouyuklieva \cite{B-LCD-40} is invalid.

The paper is organized as follows. In Section 2, we give some notations and definitions, which can be found in \cite{Huffman,MacWilliams}. In Section 3, we give a method for constructing LCD $[n-l,k-l,\geq d]$ and $[n-l,k,d-l]$ codes from a given binary $[n,k,d]$ code with $l$-dimensional hull. In addition, we give a method for constructing LCD $[n+1,k]$ and $[n,k+1]$ codes from a given binary LCD $[n,k]$ code. In Section 4, we construct binary LCD codes. In Section 5, we construct ternary LCD codes. In Section 6, we construct some new quaternary Hermitian LCD codes and the related
entanglement-assisted quantum error correction codes.
In Section 7, we conclude the paper.
All computations in this paper have been done with the computer algebra
system MAGMA \cite{magma}. We list all LCD codes with generator matrices in {\small https://docs.qq.com/doc/p/2a8a0b26ed559e57ca3f740b80d093bc40dbd544?dver=2.1.27514623}.

\section{Preliminaries}


Let $\F_q$ denote the finite field with $q$ elements, where $q$ is a prime power.
For any $\textbf x\in \F_q^N$, the Hamming weight of $\textbf x$ is the
number of nonzero components of $\textbf x$.
An $[N,K,D]$ linear code $C$ over $\F_q$ is a $K$-dimension subspace of $\F_q^N$, where $D$ is the minimum nonzero Hamming weight of $C$.
Two codes $C$ and $D$ are equivalent if there is a monomial matrix $M$ such
that $MC = D$.
The Euclidean dual code $C^{\perp_E}$ of a linear code $C$ is defined as
$$C^{\perp_E}=\{\textbf y\in \F_q^N~|~\langle \textbf x, \textbf y\rangle_E=0, {\rm for\ all}\ \textbf x\in C \},$$
where $\langle \textbf x, \textbf y\rangle_E=\sum_{i=1}^N x_iy_i$ for $\textbf x = (x_1,x_2, \ldots, x_N)$ and $\textbf y = (y_1,y_2, \ldots, y_N)\in \F_q^N$.
The Hermitian dual code $C^{\perp_H}$ of a linear code $C$ over $\F_{q^2}$ is defined as
$$C^{\perp_H}=\{\textbf y\in \F_{q^2}^N~|~\langle \textbf x, \textbf y\rangle_H=0, {\rm for\ all}\ \textbf x\in C \},$$
where $\langle \textbf x, \textbf y\rangle_H=\sum_{i=1}^N x_i\overline{y_i}$ for $\textbf x = (x_1,x_2, \ldots, x_N)$ and $\textbf y = (y_1,y_2, \ldots, y_N)\in \F_{q^2}^N$. Note that $\overline{x}=x^q$ for any $x\in \F_{q^2}$.
The Euclidean hull (resp. Hermitian hull) of the linear code $C$ is defined as
$${\rm Hull_E}(C)=C\cap C^{\perp_E}\ ({\rm resp.\ Hull_H}(C)=C\cap C^{\perp_H}).$$


If $C\cap C^{\perp_E}=\{\textbf 0\}$, the code $C$ is called Euclidean LCD.
If $C\cap C^{\perp_H}=\{\textbf 0\}$, the code $C$ is called Hermitian LCD.
These two families of codes are collectively called LCD codes.

An LCD $[N, K]$ code with the largest minimum distance among all LCD $[N, K]$ codes is  optimal LCD. An LCD $[N,K,D]$ code is called almost optimal LCD if there exists an optimal  LCD $[N,K,D+1]$ code.
Throughout this paper, let $d^E_q(n, k)$ denote the largest minimum weight among all LCD
$[n, k]$ codes over $\F_q$ $(q = 2, 3)$, and let $d_4^H(n, k)$ denote the largest minimum weight among all quaternary Hermitian LCD $[n, k]$ codes.

A vector $\textbf x = (x_1, x_2, \ldots , x_n)\in \F^n_2$
is even-like if $\sum_{i=1}^nx_i=0$ and is odd-like otherwise.
A binary code is said to be even-like if it
has only even-like codewords, and is said to be odd-like if it
is not even-like.
The following lemma comes from \cite{LCD-6} and \cite{LCD-7}.
\begin{lem}
{\rm(1)} Let $C$ be a code with the generator matrix $G$ over $\F_q$. Then $C$ is LCD
if and only if $GG^T$ is nonsingular, where $G^T$ denotes the transpose of $G$.\\
{\rm(2)} Let $C$ be a code with the generator matrix $G$ over $\F_{q^2}$. Then $C$ is Hermitian LCD if and only if $G\overline{G}^T$ is nonsingular, where $\overline{G}^T$ denotes the conjugate transpose of $G$.
\end{lem}


Let $C$ be an $[n, k, d]$ code over $\F_q$, and let $T$ be a set of $t$
coordinate positions in $C$. We puncture $C$ by deleting all the
coordinates in $T$ in each codeword of $C$. The resulting code
is still linear and has length $n - t$. We denote
the punctured code by $C^T$.
Consider the set
$C(T)$ of codewords which are $0$ on $T$; this set is a subcode of $C$. Puncturing $C(T)$ on $T$ gives
a code over $\F_q$ of length $n -t$ called the code shortened on $T$ and denoted $C_T$.

The following lemma is also valid with respect to the Hermitian inner product.


\begin{lem}\label{lem-shorten-puncture}{\rm\cite{Huffman}}
Let $C$ be an $[n,k,d]$ code over $\F_q$. Let $T$ be a set of $t$ coordinates. Then:
\begin{itemize}
  \item [(i)] $(C^\perp)_T=(C^T)^\perp$ and $(C^\perp)^T=(C_T)^\perp$, and
  \item [(ii)] if $t<d$, then $C^T$ and $(C^\perp)_T$ have dimensions $k$ and $n- t- k$, respectively.
%  \item [(iii)] if $t=d$ and $T$ is the set of coordinates where a minimum weight codeword is nonzero, then $C^T$ and $(C^\perp)_T$ have dimensions $k - 1$ and $n -d - k + 1$, respectively.
\end{itemize}
\end{lem}





\section{Methods for constructing LCD codes}

\subsection{LCD codes from shortened codes and punctured codes of linear codes}
The puncturing and shortening techniques are two very important tools for constructing new codes from old ones. In this subsection, we will use these two techniques to construct
new LCD codes with interesting and new parameters from some old linear codes.
Firstly, we prove that Lemma 22 in \cite{LCD-3} is valid with respect to the Hermitian inner product.



\begin{theorem}\label{thm-SO+LCD}
Any linear code $C$ over $\F_q$ (resp. $\F_{q^2}$) is the direct sum of a self-orthogonal code and an LCD code with respect to the Euclidean (resp. Hermitian) inner product.
\end{theorem}

\begin{proof}
We only consider the Hermitian inner product.
Let $\{\alpha_1,\ldots,\alpha_l,\alpha_{l+1},\ldots,\alpha_k\}$ be a basis of $C$ such that
$\{\alpha_1,\ldots,\alpha_l\}$ is a basis of $C_1=Hull_H(C)=C\cap C^{\perp_H}$.
Let $C_2$ is a linear code generated by $\alpha_{l+1},\alpha_{l+2},\ldots,\alpha_{k}$.
Thus $C=C_1\oplus C_2$.
For any $c\in C_2\cap C_2^{\perp_H}$, $\langle c, \alpha_i\rangle_H=0$. Implying that $c\in C^{\perp_H}$. Hence $c=0$ from $C_2\cap C^{\perp_H}=\{0\}$. Therefore, $C_2$ is LCD with respect to the Hermitian inner product.
This completes the proof.
\end{proof}

\begin{theorem}\label{LCD-(n-l,k-l)}
If there exists an $[n,k,d]$ linear code $C$ with $l$-dimensional hull. Then there exists a set of $l$ coordinates position $T$ such that the shortened $C_T$ of $C$ on $T$ is an $[n-l,k-l,\geq d]$ LCD code with respect to the Euclidean and Hermitian inner product.
\end{theorem}

\begin{proof}
Let $G$ be a generator matrix of $C$. Without loss of generality, we may assume that $$G=(I_k|A)=(e_{k,i}|a_i)_{1\leq i\leq k},$$
where $e_{k,i}$ and $a_i$ are the $i$-row of $I_k$ (the identity matrix) and $A$, respectively.
Assume that $\{r_j\}_{j=1}^l$ is a basis of $Hull(C)$ such that the first non-zero position of $r_j$ is the $i_j$-th position. Without loss of generality, we may assume that
$1\leq i_1< i_2< \cdots < i_l$. Then $i_l\leq k$; otherwise $i_l=\textbf 0$, which is a contradiction.

Let $T=\{i_1,i_2,\ldots,i_l\}$ and $J=\{1,2,\ldots,k\}\setminus T=\{j_1,j_2,\ldots,j_{k-l}\}$ such that $j_1<j_2<\cdots<j_{k-l}$.
Then we know that $\{r_j\}_{j=1}^l \cup \{(e_{k,i}|a_i)\}_{i\in J}$ is a basis of $C$.
So the code $C(T)$ by generating
$\{(e_{k,i}|a_i)\}_{i\in J}$ is an LCD code by Theorem \ref{thm-SO+LCD}.
In fact, the generator matrix for the shortened code $C_T$ on $T$ is
$$G_T=(e_{k-l,i}|a_{j_i})_{1\leq i\leq k-l},$$
 where $e_{k-l,i}$ is the $i$-row of $I_{k-l}$ and $a_{j_i}$ is the ${j_i}$-row of $A$ for $1\leq i\leq k-l$.
It follows from $C(T)$ is LCD that $C_T$ is LCD. The parameters for $C_T$ are obvious.
\end{proof}



\begin{theorem}\label{puncture}
Let $C$ be an $[n,k,d]$ linear code with $l$-dimensional hull. If $t<d$, then there exists a set of $t$ coordinate position $T$ such that the punctured code $C^T$ of $C$ on $T$ is an $[n-l,k,d^*\geq d-l]$ LCD linear code with respect to the Euclidean and Hermitian inner product.
\end{theorem}

\begin{proof}
The parameters for the punctured code $C^T$ of $C$ on $T$ are obvious from (ii) of Lemma \ref{lem-shorten-puncture}.
Since $C$ is an $[n,k]$ linear code with $l$-dimensional hull, $C^{\perp}$ is an $[n,n-k]$ linear code with $l$-dimensional hull.
According to Theorem \ref{LCD-(n-l,k-l)}, there exists a set of $l$ coordinate positions $T$ such that the shortened code $(C^\perp)_T$ of $C^\perp$ on $T$ is an LCD $[n-t,n-k-t]$ code.
It follows from (i) of Lemma \ref{lem-shorten-puncture} that $(C^T)^\perp=(C^\perp)_T$ is an LCD $[n-t,n-k-t]$ code.
It turns out that $C^T=((C^T)^\perp)^\perp$ is an LCD $[n-t,k]$ code.
\end{proof}






\begin{remark}
According to Theorem \ref{LCD-(n-l,k-l)}, $C_T$ is LCD if $T$ is an information set of $Hull(C)$.
Compared to the randomness of \cite{2-LCD-40} and \cite{H-25}, we determine the shortened set $T$.

In addition, the corollary 25 in \cite{LCD-3} proved that there exists an $[n+l,k,\geq d]$ Euclidean LCD code if there exists an $[n,k,d]$ code with $l$-dimensional Euclidean hull. Hence,
Theorem \ref{LCD-(n-l,k-l)} is more effective than \cite[Corollary 25]{LCD-3}.
\end{remark}



\subsection{Construction method 1}
In this subsection, we give a method for constructing many $[n+1,k+1]$ LCD codes with interesting parameters from a given $[n,k]$ LCD codes by modifying the method in \cite{234-lcD}.
This method can construct some LCD codes that cannot be constructed from \cite{234-lcD}.

\begin{theorem}\label{Methods}
{\rm(1)} Let $C$ be a binary LCD $[n,k]$ code with the generator matrix $G$. Let ${\bf x}\in C^\perp$. Let $C'$ be the binary linear code with the generator matrix
$$
G'=\begin{pmatrix}
1 & {\bf x} \\
 {\rm{\bf 0}} & G
\end{pmatrix}.
$$
If $wt({\bf x})$ is even, then $C'$ is a binary LCD $[n+1,k+1]$ code.

{\rm(2)} Let $C$ be a ternary LCD $[n,k]$ code with the generator matrix $G$. Let ${\bf x}\in C^\perp$. Let $C'$ be the ternary linear code with the generator matrix
$$
G'=\begin{pmatrix}
1 & {\bf x} \\
 {\rm{\bf 0}} & G
\end{pmatrix}.
$$
If $wt({\bf x})\neq 2\ ({\rm mod}\ 3)$, then $C'$ is a ternary LCD $[n+1,k+1]$ code.

{\rm(3)} Let $C$ be a quaternary Hermitian LCD $[n,k]$ code with the generator matrix $G$. Let ${\bf x}\in C^{\perp_H}$. Let $C'$ be the quaternary linear code with the generator matrix
$$
G'=\begin{pmatrix}
1 & {\bf x} \\
 {\rm{\bf 0}} & G
\end{pmatrix}.
$$
If $wt({\bf x})$ is even, then $C'$ is a quaternary Hermitian LCD $[n+1,k+1]$ code.
\end{theorem}

\begin{proof}
We only give the proof of (1), the other cases are similar.
Let $r_i$ be the $i$-th row of $G$ for $1\leq i\leq k$. Let $r'_j$ be $j$-th row of $G'$ for $1\leq j\leq k+1$.
Then we have
\begin{align*}
  \langle r'_1, r'_1\rangle_E  =&1+\langle {\bf x}, {\bf x}\rangle_E=1, \\
  \langle r'_1, r'_j\rangle_E  =&\langle {\bf x}, r_{j-1}\rangle_E=0\ {\rm for}\ 2\leq j\leq k+1,\\
  \langle r'_j, r'_j\rangle_E  =&\langle r_{j-1}, r_{j-1}\rangle_E\ {\rm for}\ 2\leq j\leq k+1,\\
  \langle r'_j, r'_{j'}\rangle_E  =&\langle r_{j-1}, r_{j'-1}\rangle_E\ {\rm for}\ 2\leq j<j'\leq k+1.
\end{align*}
Hence we obtain that
$$
G'G'^{T}=\begin{pmatrix}
1 & 0\ 0\ \cdots\ 0 \\
\textbf 0 & GG^T
\end{pmatrix}.
$$
If $C$ is a binary LCD code, then $GG^{T}$ is nonsingular. Implying that
$G'G'^{T}$ is nonsingular. Hence $C'$ is LCD. This completes the proof.
\end{proof}

It is easy to see that all binary LCD codes constructed by Theorem \ref{Methods} are odd-like.
Next, we prove that all odd-like binary LCD codes can be obtained by
the construction in Theorem \ref{Methods}.

\begin{theorem}
Any odd-like binary LCD $[n,k,d]$ code is obtained from some binary LCD $[n-1,k-1,\geq d]$ code (up to equivalence) by the construction of Theorem \ref{Methods}.
\end{theorem}

\begin{proof}
Let $C'$ be a binary LCD $[n,k,d]$ code. According to \cite[Propsition 3]{B-LCD-40}, there is at least one coordinate position $i$ such that the shortened code $C'_i$ of $C'$ on the $i$-th coordinate is a binary LCD $[n-1,k-1,\geq d]$ code. Without loss of generality, we consider that $i=1$.
Assume that $C'_1$ has the generator matrix $G'_1$. Then $C'$ has the generator matrix
$G'=\begin{pmatrix}
1 & {\bf x'} \\
\textbf 0 & G'_1
\end{pmatrix}$
for some ${\bf x'}=(x'_1,x'_2,\ldots,x'_{n-1})\in \F_2^{n-1}$.
Since $C'_1$ is a binary LCD code, $\F_2^{n-1}=C'_1\oplus {C'_1}^\perp$.
So there are ${\bf x}=(x_1,x_2,\ldots,x_{n-1})\in {C'_1}^\perp$ and ${\bf y}=(y_1,y_2,\ldots,y_{n-1})\in C'_1$ such that ${\bf x'}={\bf x}+{\bf y}$.
Let $G''=\begin{pmatrix}
1 & {\bf x} \\
\textbf 0 & G'_1
\end{pmatrix}$.
Then $G''$ is the generator matrix of $C'$.
It turns out that $wt(x)$ is even; otherwise $(1|{\bf x})\in C'\cap {C'_1}^\perp$, which is a contradiction.
This completes the proof.
\end{proof}

Finally, we give two counterexamples to show that Conjecture 1 in \cite{B-LCD-40} is invalid.

\begin{conjecture}\label{conjecture}{\rm\cite[Conjecture 1]{B-LCD-40}}
Let $k$ be an even position integer and $n>k$ be another integer. If $d_{2}^{E}(n,k)$ is even and $d_2^{E}(n-1,k)=d_2^{E}(n,k)-1$, then all binary LCD $[n,k,d_2^{E}(n,k)]$ are even-like codes.
\end{conjecture}


\begin{prop}\label{anti-example}
There exist optimal odd-like binary LCD $[14,8,4]$ and $[16,10,4]$ codes.
\end{prop}


\begin{proof}
According to \cite{HS-BLCD-1-16}, there exist optimal odd-like binary LCD $[13,7,4]$ and $[15,9,4]$ codes, which have the following generator matrices respectively
{\small\begin{center}
$\begin{pmatrix}
\setlength{\arraycolsep}{1.2pt}
    \begin{array}{ccccccccccccc}
      1&0&0&0&0&0&0&1&1&0&1&0&1 \\
        0&1&0&0&0&0&0&0&1&1&1&1&0 \\
        0&0&1&0&0&0&0&0&0&1&1&0&1 \\
        0&0&0&1&0&0&0&1&0&0&1&1&0 \\
        0&0&0&0&1&0&0&0&0&1&0&1&1 \\
        0&0&0&0&0&1&0&1&1&1&0&1&1 \\
        0&0&0&0&0&0&1&1&1&1&1&0&0 \\
    \end{array}
\end{pmatrix},~
$
$\begin{pmatrix}
\setlength{\arraycolsep}{1.2pt}
   \begin{array}{ccccccccccccccc}
    1&0&0&0&0&0&0&0&0&1&1&0&0&0&1\\
0&1&0&0&0&0&0&0&0&0&1&1&1&1&0\\
0&0&1&0&0&0&0&0&0&0&1&1&1&0&1\\
0&0&0&1&0&0&0&0&0&1&1&0&0&1&0\\
0&0&0&0&1&0&0&0&0&0&1&1&0&1&1\\
0&0&0&0&0&1&0&0&0&0&1&0&1&1&1\\
0&0&0&0&0&0&1&0&0&1&1&1&1&1&1\\
0&0&0&0&0&0&0&1&0&1&1&0&1&0&0\\
0&0&0&0&0&0&0&0&1&1&1&1&0&0&0\\
   \end{array}
\end{pmatrix}.
$
\end{center}}
Applying Theorem \ref{Methods} to the LCD $[13,7,4]$ code, a binary LCD $[14,8,4]$ code $C'$ is constructed, where ${\bf x}=(1 0 0 1 1 1 0 0 0 1 1 0 0)$. The code $C'$ has weight distribution as:
$$[ \langle0, 1\rangle, \langle4, 24\rangle, \langle5, 36\rangle, \langle6, 36\rangle, \langle7, 60\rangle, \langle8, 45\rangle, \langle9, 28\rangle, \langle10, 20\rangle, \langle11,4\rangle, \langle12, 2\rangle].$$
Similarly, a binary LCD $[16,10,4]$ code is constructed, where ${\bf x}=(1 1 1 1 1 1 0 1 1 0 0 1 1 1 1)$, which has weight distribution as:
\begin{center}
$[ \langle0, 1\rangle, \langle4, 43\rangle, \langle5, 81\rangle, \langle6, 96\rangle, \langle7, 189\rangle, \langle8, 207\rangle, \langle9, 162\rangle,$\\ $\langle10, 144\rangle,
\langle11, 66\rangle, \langle12, 21\rangle, \langle13, 13\rangle, \langle15, 1\rangle ].$
\end{center}
In the sense of equivalence, these two LCD code are also constructed by \cite{HS-BLCD-1-16}. Here, we just show that we can get such codes by our method.
\end{proof}




\begin{remark}
Let $n=14,k=8$.
According to \cite{HS-BLCD-1-16}, $d_{2}^E(14,8)=4$ and $d_{2}^E(13,8)=3$. The condition of conjecture \ref{conjecture} is satisfied. But the binary LCD $[14,8,4]$ code in Proposition \ref{anti-example} is odd-like. Therefore, Conjecture {\rm\ref{conjecture}} is invalid.
Similarly, the odd-like binary LCD $[16,10,4]$ is also a counterexample of Conjecture \ref{conjecture}.
\end{remark}


\subsection{Construction method 2}
In this subsection, we give a method for constructing many $[n,k+1]$ LCD codes with interesting parameters from a given $[n,k]$ LCD codes by modifying the method in \cite{34-LCD}.
\begin{theorem}\label{Methods-2}
{\rm(1)} Let $C$ be a binary LCD $[n,k]$ code with the generator matrix $G$. Let ${\bf y}\in C^\perp$. Let $C'$ be the binary linear code with the generator matrix
$$
G'=\begin{pmatrix}
 {\bf y} \\
  G
\end{pmatrix}.
$$
If $wt({\bf y})$ is odd, then $C'$ is a binary LCD $[n,k+1]$ code.

{\rm(2)} Let $C$ be a ternary LCD $[n,k]$ code with the generator matrix $G$. Let ${\bf y}\in C^\perp$. Let $C'$ be the ternary linear code with the generator matrix
$$
G'=\begin{pmatrix}
 {\bf y} \\
  G
\end{pmatrix}.
$$
If $wt({\bf y})\neq 0\ ({\rm mod}\ 3)$, then $C'$ is a ternary LCD $[n,k+1]$ code.

{\rm(3)} Let $C$ be a quaternary Hermitian LCD $[n,k]$ code with the generator matrix $G$. Let ${\bf y}\in C^{\perp_H}$. Let $C'$ be the quaternary linear code with the generator matrix
$$
G'=\begin{pmatrix}
 {\bf y} \\
 G
\end{pmatrix}.
$$
If $wt({\bf y})$ is odd, then $C'$ is a quaternary Hermitian LCD $[n,k+1]$ code.
\end{theorem}

\begin{proof}
The proof is similar to that of Theorem \ref{Methods}, so we omit it here.
\end{proof}

Next, we prove that most LCD codes can be obtained by
the construction in Theorem \ref{Methods-2}.

\begin{theorem}\label{up-2}
Any binary odd-like LCD $[n,k,d]$ code is obtained from some binary LCD $[n,k-1,\geq d]$ code (up to equivalence) by the construction of Theorem \ref{Methods-2}.
\end{theorem}

\begin{proof}
Let $C'$ be a binary odd-like LCD $[n,k,d]$ code. From \cite[Theorem 3]{LCD-new}, there exists a basis $c_1,c_2,\ldots,c_{k}$ of $C'$ such that $c_i \cdot c_j=1$ if $i=j$ and $c_i\cdot c_j=0$ otherwise. Let $C$ be a binary code with the generator matrix $G=(c_{i+1})_{1\leq i\leq k-1}$, where $c_{i+1}$ is the $i$-row of $G$. From \cite[Theorem 3]{LCD-new}, $C$ is LCD. Implying that $\F_{2}^n=C \oplus C^\perp$.
So there are ${\bf x}=(x_1,x_2,\ldots,x_n)\in C$ and ${\bf y}=(y_1,y_2,\ldots,y_n)\in C^\perp$ such that $c_1={\bf x}+{\bf y}$.
Let $G'=\begin{pmatrix}
{\bf y} \\
 G
\end{pmatrix}$.
Then $G'$ is the generator matrix of $C$.
It turns out that $wt({\bf y})$ is odd; otherwise ${\bf y}\in C'\cap C'^{\perp}$, which is a contradiction.
This completes the proof.
\end{proof}

\begin{theorem}
Any ternary Euclidean (resp. quaternary Hermitian) LCD $[n,k]$ code is obtained from some ternary Euclidean (resp. quaternary Hermitian) LCD $[n,k-1]$ code by the construction of Theorem \ref{Methods-2}.
\end{theorem}

\begin{proof}
From \cite[Propsition 4]{34-LCD}, we know that any ternary (resp. quaternary Hermitian) LCD $[n,k]$ code contains an LCD $[n,k-1]$ subcode.
The rest of the proof is similar to that of Theorem \ref{up-2}, so we omit it here.
\end{proof}






\section{New binary LCD codes}
\subsection{Some important inequalities}
Let $d^E_2 (n, k)$ denote the largest minimum weight among all binary Euclidean LCD $[n, k]$ codes.
By adding zero column, the following inequality is obvious.

\begin{lem}\label{lem-inequality-1}
 Suppose that $ k\leq n$. Then $d^E_2 (n+1, k)\geq d^E_2 (n, k)$.
\end{lem}


The following lemma was proved in \cite[Theorem 8]{LCD-new}.
\begin{lem}{\rm\cite{LCD-new}}\label{lemma-leq}
Suppose that $2\leq k\leq n$. Then
$d^E_2(n,k)\leq d^E_2(n,k-1).$
\end{lem}


The following inequalities are also important.
\begin{lem}{\rm\cite{B-LCD-40}}\label{lemma-leq-2}
Let $k$ and $n$ be two integers such that $1\leq k\leq n$. Then we have
\begin{itemize}
  \item [(1)] If $k$ is odd, then $d^E_2(n,k)\leq d^E_2(n-1,k-1).$
  \item [(2)] If $k$ is even and $d^E_2(n,k)$ is odd, then $d^E_2(n+1,k)\geq d^E_2(n,k)+1$.
  \item [(3)] If $d^E_2(n,k)$ is odd, then $d^E_2(n+2,k)\geq d^E_2(n,k)+1$.
\end{itemize}
\end{lem}



\begin{cor}\label{cor-4.4}
Let $k$ and $n$ be two integers such that $2\leq k\leq n$. Then we have
$$d^H_2(n,k)\leq \max\{d^H_2(n-1,k-1),d^H_2(n-2,k-2)\}.$$
\end{cor}
\begin{proof}
Assume that $C$ is a binary $[n,k,d]$ LCD code. Let $C_i$ be the
shortened code of $C$ on $i$-th coordinate for some $i\in \{1,2,\ldots,n\}$. By \cite[Proposition 5]{B-LCD-40}, we know that $\dim(Hull(C_i))\leq 1$.
If $\dim(Hull(C_i))=0$, then there exists a binary $[n-1,k-1,d]$ LCD code.
If $\dim(Hull(C_i))=1$, then there exists a binary $[n-1,k-1,d]$ code with one-dimension hull. By Theorem \ref{LCD-(n-l,k-l)}, there exists a binary $[n-2,k-2,d]$ LCD code.
 This completes the proof.
\end{proof}


Let $C^{l}_{[n,k,d]_2}$ denote a binary $[n,k,d]$ linear code with the generator matrix $G^{l}_{[n,k,d]_2}$, where $l$ is the dimension of hull for such code, when $l=0$, let $C_{[n,k,d]_2}=C^{0}_{[n,k,d]_2}$ denote a binary LCD $[n,k,d]$ code.

\subsection{New binary LCD codes}


From \cite{B-LCD-40}, only the exact value of $d_{2}^E(29,11)$ remains unknown for $n=29$.


\begin{prop}\label{29}
There is a binary LCD $[29,11,9]$ code.
\end{prop}
\begin{proof}
By the MAGMA function BKLC, there is a binary $[30,11,10]$ code $C^{1}_{[30,11,10]_2}$ with one-dimensional hull. The shortened code $(C^{1}_{[30,11,10]_2})_{\{1\}}$ on the first coordinate position is a binary LCD $[29,10,10]$ code $C_{[29,10,10]_2}$.
Applying Theorem \ref{Methods-2} to the code $C_{[29,10,10]_2}$, a binary LCD $[29,11,9]$ code $C_{[29,11,9]_2}$ is constructed by adding ${\bf y}$ to the code $C_{[29,10,10]_2}$, where
${\bf y}=(0 0 0 1 0 1 0 1 0 1 1 1 1 0 1 0 1 0 0 1 1 0 0 0 1 0 0 0 1).$
\end{proof}

\begin{prop}\label{[30,12,9]-[30,15,7]}
There are binary LCD $[30,11,9]$, $[30,15,7]$ and $[33,11,11]$ codes.
\end{prop}

\begin{proof}
It follows from Lemma \ref{lem-inequality-1} and Proposition \ref{29} that there is a binary LCD $[30,11,9]$ code.
By the MAGMA function BKLC, there is a binary $[36,21,7]$ code $C_{[36,21,7]_2}^6$ with 6-dimensional hull.
Applying Theorem \ref{LCD-(n-l,k-l)} to the code $C_{[36,21,7]_2}^6$, there is a binary LCD $[30,15,7]$ code $C_{[30,15,7]_2}$. In fact, we may assume that $C_{[30,15,7]_2}$ is the shortened code $(C_{[36,21,7]_2}^6)_T$ of $C_{[36,21,7]_2}^6$, where $T=\{1,2,3,5,16,18\}$.

According to \cite{2-LCD-40}, there is a binary LCD $[34,12,11]$ code $C_{[34,12,11]_2}$.
So $C_{[34,12,11]_2}$ is an odd-like binary LCD code. By \cite[Propsition 3]{2-LCD-40}, there is at least one coordinate position $i$ such that the shortened code $(C_{[34,12,11]_2})_{\{i\}}$ of $C_{[34,12,11]_2}$ on the $i$-th coordinate is a binary LCD $[33,11,\geq 11]$ code.
\end{proof}




{\small
\begin{center}
 \begin{threeparttable}
\begin{tabular}{ccc}
\multicolumn{3}{c}{{\rm Table 1: The binary LCD codes from Theorem \ref{LCD-(n-l,k-l)}}}\\
\hline
    \makebox[0.3\textwidth][c]{{\rm The\ linear\ code}\ $C_{[n,k,d]_2}^l$}& \makebox[0.3\textwidth][c]{The set $T$} &
    \makebox[0.3\textwidth][c]{The parameters for $C_T$} \\
    \hline\hline
$C^{1}_{[30,11,10]_2} $&$ \{1\} $&$ [29,10,10]$\\

$C_{[36,21,7]_2}^6 $&$ \{1,2,3,5,16,18\} $&$ {\bf[30,15,7]}$\\

$C^{1}_{[32,21,6]_2} $&$ \{1\} $&$ [31,20,6]$\\

$C^{0}_{[33,22,6]_2} $&$\backslash $&$ [33,22,6]$\\

$C^{1}_{[34,15,9]_2} $&$ \{1\} $&$ [33,14,9]$\\

$C^{0}_{[34,22,6]_2} $&$ \backslash $&$ [34,22,6]$\\

$C^{7}_{[43,27,7]_2} $&$ \{1,2,3,5,9,17,25\} $&$ {\bf[36,20,7]}$\\

$C^{3}_{[40,13,13]_2} $&$ \{2,3,5\} $&$ {\bf[37,10,13]}$\\

$C^{8}_{[45,29,7]_2} $&$ \{1,2,3,5,9,16,20,23\} $&$ {\bf[37,21,7]}$\\
\hline
\end{tabular}
\begin{tablenotes}
\footnotesize
\item [1] {\small We obtain some binary LCD codes by Theorem \ref{LCD-(n-l,k-l)}. All results are listed in Table 1, where the linear code $C_{[n,k,d]_2}^l$ is the best-known binary $[n,k,d]$ code from MAGMA \cite{magma}. It is worth noting that the code in MAGMA may not be in standard form, we need to use the MAGMA function StandardForm to act on it.}
\end{tablenotes}
\end{threeparttable}
\end{center}

}

\begin{remark}
Although the code $C_{[29,11,9]_2}$ has the same parameters with the $[29,11,9]$ code in Proposition 3 of \cite{2-LCD-40}, their construction methods are different. In addition, they are not equivalent by Magma \cite{magma}. It can be seen that the code $C_{[30,11,9]_2}$ and the $[30,11,9]$ code in Proposition 4 of \cite{2-LCD-40} are also not equivalent. Hence we can consider that the binary LCD $[29,11,9]$ and $[30,11,9]$ codes we obtained are new.
Similarly, we can consider that the binary LCD codes are in Tables 1-3 new.
\end{remark}





{\small
\begin{center}
 \begin{threeparttable}
\begin{tabular}{cccc}
\multicolumn{4}{c}{{\rm Table 2: The binary LCD codes from Theorem \ref{Methods} and Theorem \ref{Methods-2}}}\\
\hline
    \makebox[0.15\textwidth][c]{{\rm The\ given}\ $C_{[n,k,d]_2}$}& \makebox[0.22\textwidth][c]{The vector ${\bf x}$ or ${\bf y}$} &
    \makebox[0.12\textwidth][c]{$C'$} &
    \makebox[0.12\textwidth][c]{References}\\
    \hline\hline
[30,15,7] & (110001101111101110111111001001) & {\bf[31,16,7]} & {\rm Theorem}\ \ref{Methods} \\

[31,16,7] & (1110101000110101100110111011001) & ${\bf[32,17,7]}$& {\rm Theorem}\ \ref{Methods} \\

[32,17,7] & (11010011101011010101111000110111) & {\bf[33,18,7]}& {\rm Theorem}\ \ref{Methods} \\

[33,14,9] & (111000000011111101010110010000110) & [34,15,9]& {\rm Theorem}\ \ref{Methods} \\

[34,18,8] & (1100000000000100001100101101010110) & {\bf[35,19,7]}&{\rm Theorem}\ \ref{Methods} \\

[34,22,6] &  (1111111111111111111111111111111111) & {\bf[35,23,6]}& {\rm Theorem}\ \ref{Methods} \\

[36,16,10] &  (110100000100000010011010100001101110) & [37,17,9]&{\rm Theorem}\ \ref{Methods} \\

[37,21,7] & (0011101010010010110110100111111001111) & {\bf[38,22,7]}&{\rm Theorem}\ \ref{Methods} \\

[38,10,14] &  (11010011110000000000000000000011101100) & [39,11,13]&{\rm Theorem}\ \ref{Methods} \\

[38,22,7] & (10110110101011110010110101001001010100) & {\bf[39,23,7]}&{\rm Theorem}\ \ref{Methods} \\

[29,10,10] & (00010101011110101001100010001) & [29,11,9] &{\rm Theorem}\ \ref{Methods-2} \\

[31,20,6] & (0000111101010000001110010011100) & [31,21,5] &{\rm Theorem}\ \ref{Methods-2} \\
\hline
\end{tabular}
\begin{tablenotes}
\footnotesize
\item [2] {\small We obtain some new binary LCD codes by construction methods in Theorems \ref{Methods} and \ref{Methods-2}. All results are listed in Table 2, where the given code $C_{[n,k,d]_2}$ in Table 2 is from Tables 1-2.}
\end{tablenotes}
\end{threeparttable}
\end{center}}








{\small
\begin{center}
 \begin{threeparttable}
\begin{tabular}{ccc||ccc}
\multicolumn{6}{c}{{\rm Table 3: The binary LCD codes from some inequalities}}\\
\hline
    \makebox[0.12\textwidth][c]{{\rm The\ given}\ $C$}&
    \makebox[0.12\textwidth][c]{$C'$} &
    \makebox[0.1\textwidth][c]{References} &
   \makebox[0.12\textwidth][c]{{\rm The\ given}\ $C$}&
    \makebox[0.12\textwidth][c]{$C'$} &
    \makebox[0.1\textwidth][c]{References} \\
    \hline\hline
[29,11,9] & [30,11,9] & {\rm Lemma\ \ref{lem-inequality-1}}&
[35,21,6] & [36,21,6] & {\rm Lemma\ \ref{lem-inequality-1}}\\

[29,11,9] & [31,11,10] & {\rm (3)\ of\ Lemma\ \ref{lemma-leq-2}}&[35,23,6] & [36,23,6] & {\rm Lemma\ \ref{lem-inequality-1}}\\

[31,16,7] & [31,15,7] & {\rm Lemma\ \ref{lemma-leq}}&[36,18,8]&[37,18,8]&{\rm Lemma\ \ref{lem-inequality-1}}\\

[31,16,7] & {\bf[32,16,8]} & {\rm (2)\ of\ Lemma\ \ref{lemma-leq-2}}&
[35,19,7]&{\bf[37,19,8]}&{\rm (3)\ of\ Lemma\ \ref{lemma-leq-2}}\\

[32,16,8] & {\bf[32,15,8]} & {\rm Lemma\ \ref{lemma-leq}}&
[36,20,7]&{\bf[37,20,8]}& {\rm (2)\ of\ Lemma\ \ref{lemma-leq-2}}\\

[31,21,5] & [32,21,5] & {\rm Lemma\ \ref{lem-inequality-1}}&
[36,23,6]&[37,23,6]&{\rm Lemma\ \ref{lem-inequality-1}}\\

[32,15,8] & [33,15,8] & {\rm Lemma\ \ref{lem-inequality-1}}&
[37,10,13]&{\bf[38,10,14]}&{\rm (2)\ of\ Lemma\ \ref{lemma-leq-2}}\\

[32,16,8] & {\bf[33,16,8]} & {\rm Lemma\ \ref{lem-inequality-1}}&
[37,17,9]&[38,17,9]&{\rm Lemma\ \ref{lem-inequality-1}}\\

[32,17,7] & {\bf[33,17,7]} & {\rm Lemma\ \ref{lem-inequality-1}}&
[37,19,8]&[38,19,8]&{\rm Lemma\ \ref{lem-inequality-1}}\\

[33,22,6] & {\bf[33,21,6]} & {\rm Lemma\ \ref{lemma-leq}}&[37,20,8]&{\bf[38,20,8]}&{\rm Lemma\ \ref{lem-inequality-1}}\\

[33,18,7] & {\bf[34,18,8]} & {\rm (2)\ of\ Lemma\ \ref{lemma-leq-2}}&
[37,21,7]&[38,21,7]&{\rm Lemma\ \ref{lem-inequality-1}}\\

[34,18,8] & {\bf[34,17,8]} & {\rm Lemma\ \ref{lemma-leq}}&[37,23,6]&[38,23,6]&{\rm Lemma\ \ref{lem-inequality-1}}\\

[34,17,8] & [34,16,8] & {\rm Lemma\ \ref{lemma-leq}}&
[38,10,14]&[39,10,14]&{\rm Lemma\ \ref{lem-inequality-1}}\\

[33,21,6] & [34,21,6] & {\rm Lemma\ \ref{lem-inequality-1}}&
[38,20,8]&[39,20,8]&{\rm Lemma\ \ref{lem-inequality-1}}\\

[34,17,8] & {\bf[35,17,8]} & {\rm Lemma\ \ref{lem-inequality-1}}&
[38,22,7]&{\bf[39,22,8]}&{\rm (2)\ of\ Lemma\ \ref{lemma-leq-2}}\\

[34,18,8] & {\bf[35,18,8]} & {\rm Lemma\ \ref{lem-inequality-1}}&[39,22,8]&{\bf[39,21,8]}&{\rm Lemma\ \ref{lemma-leq}}\\

[34,21,6] & {\bf[35,21,6]} & {\rm Lemma\ \ref{lem-inequality-1}}&
[39,10,14]&[40,10,14]&{\rm Lemma\ \ref{lem-inequality-1}}\\

[34,15,9] & [36,15,10] & {\rm (3)\ of\ Lemma\ \ref{lemma-leq-2}}&[39,11,13]&[40,11,13]&{\rm Lemma\ \ref{lem-inequality-1}}\\

[35,17,8] & [36,17,8] & {\rm Lemma\ \ref{lem-inequality-1}}&
[39,21,8]&{\bf[40,21,8]}&{\rm Lemma\ \ref{lem-inequality-1}}\\

[35,18,8] & {\bf[36,18,8]} & {\rm Lemma\ \ref{lem-inequality-1}}&[39,22,8]&{\bf[40,22,8]} &{\rm Lemma\ \ref{lem-inequality-1}}\\

[35,19,7] & [36,19,7] & {\rm Lemma\ \ref{lem-inequality-1}}&
[39,23,7]&[40,23,7]&{\rm Lemma\ \ref{lem-inequality-1}}\\
\hline
\end{tabular}
\begin{tablenotes}
\footnotesize
\item [3] {\small We use some inequalities to obtain some new binary LCD codes, all the results are listed in Table 3, where the given code $C$ in Table 3 is from Tables 1-3.}
\end{tablenotes}
\end{threeparttable}
\end{center}}









\begin{remark}
Some codes in Tables 1-3 are optimal according to \cite{codetables}.
For example, the best-known binary linear code of the length 39 with the dimension 23 in \cite{codetables} has the minimum distance 7.
The binary LCD code we constructed also has the same parameters.
\end{remark}




{\small
\begin{center}
\begin{tabular}{cccccccccc}
\multicolumn{10}{c}{{\rm Table 4: Bounds on the minimum diatance of binary LCD codes}}\\
\hline
    \makebox[0.07\textwidth][c]{$n\backslash k$}& \makebox[0.07\textwidth][c]{9} & \makebox[0.07\textwidth][c]{10}& \makebox[0.07\textwidth][c]{11}& \makebox[0.07\textwidth][c]{12}& \makebox[0.07\textwidth][c]{13}& \makebox[0.07\textwidth][c]{14}& \makebox[0.07\textwidth][c]{15}& \makebox[0.07\textwidth][c]{16}& \makebox[0.07\textwidth][c]{17}\\
    \hline\hline
29   & 10& 10 &${\bf9^\diamond}$ & 8& 8  & 8  & 6  &  6 & 6\\

30   & 11& 10 &${\bf9^\diamond}$-10 & 9& 8  & 8  & ${\bf7^*} $ &  6 & 6\\

31   & 11& 10 &  ${\bf10^\diamond}$ &10&9&8 & ${\bf7^\diamond}$-8& ${\bf 7^*}$ & 6 \\

32   & 12& 11 &  10  & 10&9-10&8-9 & ${\bf8^*}$& ${\bf8^*}$ & ${\bf 7^*}$\\

33  & 12& 12 &  ${\bf11^*}$  & 10&10& 9-10& ${\bf8^\diamond}$-9& ${\bf 8^*} $& ${\bf 7^*}$-8\\

34   & 13& 12 &  11-12  & 11&10& 10&  ${\bf9^\diamond}$-10& ${\bf8^\diamond}$-9 & ${\bf8^*}$\\

35   & 13-14& 12-13 &  12  & 12&10-11& 10& 9-10& 9-10 & ${\bf8^*}$\\

36   & 14& 12-14 &  12-13  & 12&10-12& 10-11& ${\bf10^\diamond}$& 10 & ${\bf8^\diamond}$-9\\

37   & 14& ${\bf13^*}$-14 &  12-14& 12-13&10-12& 10-12&10-11&10 & ${\bf9^\diamond}$-10\\

38  & 14-15& ${\bf14^*}$ &  12-14& 12-14&11-12& 10-12&10-12&10-11 &${\bf9^\diamond}$-10\\

39   & 14-16& ${\bf14^\diamond}$-15 &  ${\bf13^\diamond}$-14& 12-14&11-13& 11-12&10-12&10-12 & 10-11\\

40   & 15-16& ${\bf14^\diamond}$-16 &  ${\bf13^\diamond}$-15& 13-14&12-14& 11-13&10-12&10-12 & 10-12\\
\hline
\end{tabular}
\end{center}

}

{\small
\begin{center}
 \begin{threeparttable}
\begin{tabular}{cccccccccccccc}
\multicolumn{14}{c}{{\rm Table 5: Bounds on the minimum diatance of binary LCD codes}}\\
\hline
    \makebox[0.05\textwidth][c]{$n\backslash k$}& \makebox[0.05\textwidth][c]{18} & \makebox[0.05\textwidth][c]{19}& \makebox[0.05\textwidth][c]{20}& \makebox[0.05\textwidth][c]{21}& \makebox[0.05\textwidth][c]{22}& \makebox[0.05\textwidth][c]{23}& \makebox[0.05\textwidth][c]{24}& \makebox[0.03\textwidth][c]{25}& \makebox[0.03\textwidth][c]{26}& \makebox[0.03\textwidth][c]{27}& \makebox[0.03\textwidth][c]{28}& \makebox[0.03\textwidth][c]{29}& \makebox[0.03\textwidth][c]{30}\\
    \hline\hline
29   & 6  & 5& 4  & 4  & 4  & 3&     2    & 2  & 2  & 2  & 2  &  1 & \\

30   & 6  & 5-6& 5  & 4  & 4  & 4 &     3     & 2  & 2  & 2  & 2  &  1 & 1\\

31   & 6& 6& 6& ${\bf5^\diamond}$ & 4  & 4  &  4    &   3   &2&2 & 2& 2 & 2 \\

32   & 6& 6& 6&  ${\bf5^\diamond}$-6& 5& 4 &  4     & 3-4&3&2 & 2& 2 & 2\\

33   & ${\bf 7^*}$& 6& 6&  ${\bf6^*}$&  6&  5 & 4  & 4&4& 3& 2& 2 & 2\\

34   & ${\bf 8^*}$& 6-7& 6& ${\bf6^\diamond}$&  6&  5-6 &  4  & 4&4& 3-4& 3&2 & 2\\

35   & ${\bf 8^*}$& ${\bf 7^*}$-8& 6-7& ${\bf 6^*}$&  6& ${\bf 6^*} $&   5  & 4&4& 4& 4& 3 & 2\\

36   &  ${\bf8^*}$& ${\bf7^\diamond}$-8& ${\bf 7^*}$-8&  ${\bf6^\diamond}$-7&  6&  ${\bf6^\diamond}$ &  6  &  5&4& 4& 4& 3-4 & 3\\

37   &  ${\bf8^\diamond}$-9& ${\bf 8^*}$& ${\bf 8^*}$& ${\bf 7^*}$-8& 6-7&  ${\bf6^\diamond}$ & 6&  5-6& 5& 4&4&4 & 4\\

38   & 9-10&  ${\bf8^\diamond}$-9& $ {\bf8^*}$& $ {\bf7^\diamond}$-8& ${\bf7^*}$-8& ${\bf6^\diamond}$-7 &   6& 6& 6&  5&4&4 & 4\\

39   & 10&  9-10& ${\bf8^\diamond}$-9& ${\bf 8^*}$& ${\bf8^*}$& ${\bf7^*}$-8 &   6-7&  6&6&  5-6& 5&4 & 4\\

40   & 10-11& 9-10&  9-10& ${\bf 8^*}$-9&${\bf8^*}$ & ${\bf7^\diamond}$-8 &  6-8&  6-7 & 6 &  6& 6&5 & 4\\
\hline
\end{tabular}
\begin{tablenotes}
\footnotesize
\item [4-5] {\small We give Tables 4 and 5
 by combining Tables 1-3, \cite[Tables 1-3]{2-LCD-40} and \cite[Tables 1-2]{B-LCD-40}.
Furthermore, the diamond ``$\diamond$'' indicates that the corresponding binary LCD code is
inequivalent to the code in \cite{2-LCD-40} even though they have the same parameters.
The asterisk ``$*$'' indicates that the corresponding binary LCD code has new parameters comparing with \cite{B-LCD-40} and \cite{2-LCD-40}.
For example, the best-known binary LCD code of the length 34 with the dimension 18 in \cite{2-LCD-40} has the minimum distance 6,
while the minimum distance of binary LCD codes with the same length and dimension we constructed is 8.}
\end{tablenotes}
\end{threeparttable}
\end{center}

}












\section{New ternary LCD codes}
\subsection{Some important inequalities}
Let $d^E_3(n,k)$ denote the largest minimum weight among all ternary Euclidean LCD $[n, k]$ codes. By adding zero column, we have $d^E_3 (n+1, k)\geq d^E_3 (n, k)$.

\begin{lem}{\rm\cite{34-LCD}}\label{lemma-leq-ternary}
Suppose that $2\leq k\leq n$. Then
$d^E_3(n,k)\leq d^E_3(n,k-1).$
\end{lem}

\begin{lem}{\rm\cite{ter-11-19}}
If $20\leq n\leq 25$, then we have
$$d_3^E(n,n-2)=2,\ d_3^E(n,n-3)=2,\ d_3^E(n,n-4)=3.$$
\end{lem}

According to \cite{AH-TLCD-1-10}, we know that
$$d_3^E(n,n-1)=\left\{
  \begin{array}{cc}
    1,& {\rm if}\ n\equiv 0\ ({\rm mod}\ 3), \\
    2,& {\rm otherwise.} \\
  \end{array}
\right.$$

\begin{prop}\label{prop-ternary}
Let $C$ be a ternary linear $[n,k,d]$ code and $\dim(C\cap C^\perp)=s$. If $C_i$ is the shortened code of $C$ in $i$-th coordinate for some $1\in \{1,\ldots,n\}$, then $dim(C_i\cap C_i^\perp)\leq s+1$.
\end{prop}

\begin{proof}
Without loss of generality, we can consider the shortened code $C_1$ of $C$ on the first coordinate. If all codewords have $0$ as a first coordinate, then $\dim(C_1\cap C_1)=\dim(C\cap C^\perp)=s$.
Otherwise $C=(0|C_1)\cup(1|u+C_1)\cup(2|2u+C)$ for a codeword $(1,u)\in C$.
Let $\mathcal{H}=C\cap C^\perp$ and $\mathcal{H}_1=C_1\cap C_1^\perp$.
There are two possibilities for $\mathcal{H}$, namely $\mathcal{H}=(0|\mathcal{H}')$ or $\mathcal{H}=(0|\mathcal{H}')\cup(1|v+\mathcal{H}')\cup(2|2v+\mathcal{H}')$.
In both cases, $\mathcal{H}'\subseteq\mathcal{H}_1$.
\begin{itemize}
  \item If $\mathcal{H}'=\mathcal{H}_1$, then $\dim(\mathcal{H}_1)=\dim(\mathcal{H})$ or $\dim(\mathcal{H})-1$.
  \item If $\mathcal{H}'\subsetneqq\mathcal{H}_1$. Take $y_1,y_2\in \mathcal{H}_1\backslash \mathcal{H}'\subseteq C_1\cap C_1^\perp$.
      It follows from $y_i\in C_1$ that $(0,y_i)\in C$.
      If follows from $y_i\in C_1^\perp=(C^\perp)^1$ that $(\lambda_i,y_i)\in C^\perp$ for some $\lambda_i\in \F_3$.
      If $\lambda_i=0$, then $(0,y_i)\in C\cap C^\perp=\mathcal{H}$. Implying that $y_i\in \mathcal{H}'$, which is a contradiction. Hence $\lambda_i\neq 0$. So $(0,\lambda_2y_1-\lambda_1y_2)\in \mathcal{H}$, which implies that $\lambda_2y_1-\lambda_1y_2\in \mathcal{H}'$. It turns out that $\mathcal{H}_1=\mathcal{H}'\cup(y_1+\mathcal{H}')\cup(2y_1+\mathcal{H}')$ and $\dim(\mathcal{H}_1)=\dim(\mathcal{H}')+1$.
      Since $\dim(\mathcal{H}')=\dim(\mathcal{H})$ or $\dim(\mathcal{H})-1$, we have $\dim(\mathcal{H}_1)=\dim(\mathcal{H})$ or $\dim(\mathcal{H})+1$.
\end{itemize}
 This completes the proof.
\end{proof}



\begin{cor}\label{3-cor-leq}
Let $k$ and $n$ be two integers such that $2\leq k\leq n$. Then we have
$$d^E_3(n,k)\leq \max\{d^E_3(n-1,k-1),d^E_3(n-2,k-2)\}.$$
\end{cor}
\begin{proof}
The proof is similar to that of Corollary \ref{cor-4.4}. The main difference is that we use Proposition \ref{prop-ternary} instead of \cite[Proposition 5]{B-LCD-40}.
\end{proof}

\subsection{New ternary LCD codes}

Firstly, we give some known ternary LCD codes from \cite{ter-11-19}.
\begin{center}
$13\leq d_3(23,4)\leq 14$, $d_3(24,4)=15$, $15\leq d_3(25,4)\leq 16$.
\end{center}


According to \cite{ter-11-19}, there are ternary LCD codes $C_{[19,6,9]_3}$, $C_{[20,5,11]_3}$, $C_{[20,6,10]_3}$ and $C_{[20,8,8]_3}$, they have the following generator matrices respectively.
\begin{center}
{\small
$
\begin{pmatrix}
\setlength{\arraycolsep}{1.2pt}
\begin{array}{ccccccccccccccccccc}
 1&0&0&0&0&0&0&0&0&0&0&1&1&1&1&1&1&1&1\\
0&1&0&0&0&0&0&0&0&1&1&2&2&2&1&1&1&0&0\\
0&0&1&0&0&0&1&1&1&2&1&1&1&1&0&0&0&0&0\\
0&0&0&1&0&0&1&2&2&2&1&2&0&0&1&1&0&0&0\\
0&0&0&0&1&0&1&2&0&1&2&1&1&0&1&0&0&1&0\\
0&0&0&0&0&1&1&1&2&1&0&1&0&1&2&0&1&1&0
\end{array}
\end{pmatrix}
$,\
$
\begin{pmatrix}
\setlength{\arraycolsep}{1.2pt}
\begin{array}{cccccccccccccccccccc}
1&0&0&0&0&0&0&0&0&0&1&1&1&1&1&1&1&1&1&1\\
0&1&0&0&0&0&0&1&1&1&0&0&0&1&1&1&1&2&2&2\\
0&0&1&0&0&0&1&2&2&1&2&2&1&0&2&1&1&2&2&1\\
0&0&0&1&0&1&1&2&1&1&0&2&2&2&2&1&0&0&0&0\\
0&0&0&0&1&1&2&0&1&2&2&0&2&1&2&1&2&1&0&1
\end{array}
\end{pmatrix}
$,}
\end{center}
\begin{center}
{\small$
\begin{pmatrix}
\setlength{\arraycolsep}{1.2pt}
\begin{array}{cccccccccccccccccccc}
1&0&0&0&0&0&0&0&0&0&0&1&1&1&1&1&1&1&1&1\\
0&1&0&0&0&0&0&0&1&1&1&2&2&1&1&1&1&0&0&0\\
0&0&1&0&0&0&1&1&2&1&1&1&1&1&0&0&0&1&0&0\\
0&0&0&1&0&0&1&2&0&2&1&1&1&0&1&1&0&0&1&0\\
0&0&0&0&1&0&0&1&1&1&2&1&0&0&2&0&1&1&1&0\\
0&0&0&0&0&1&1&2&2&0&0&1&2&2&1&2&0&0&0&1
\end{array}
\end{pmatrix},\
\begin{pmatrix}
\setlength{\arraycolsep}{1.2pt}
\begin{array}{cccccccccccccccccccc}
1&0&0&0&0&0&0&0&0&0&0&0&0&1&1&1&1&1&1&1\\
0&1&0&0&0&0&0&0&0&0&1&1&1&1&1&1&1&0&0&0\\
0&0&1&0&0&0&0&0&1&1&1&2&2&1&1&0&0&0&0&0\\
0&0&0&1&0&0&0&0&1&2&2&1&0&1&0&2&1&0&0&0\\
0&0&0&0&1&0&0&0&1&2&2&2&1&0&0&1&0&1&0&0\\
0&0&0&0&0&1&0&0&1&1&2&0&2&2&0&2&0&0&2&0\\
0&0&0&0&0&0&1&0&1&2&0&1&2&2&0&0&0&1&1&0\\
0&0&0&0&0&0&0&1&0&2&0&2&0&2&2&2&1&1&1&0
\end{array}
\end{pmatrix}.
$
}
\end{center}








\begin{prop}
There are ternary LCD $[20,7,9]$ and $[20,12,6]$ codes.
\end{prop}

\begin{proof}
Applying Theorem \ref{Methods} to the generator matrix $G_{[19,6,9]_3}$ of the ternary LCD code $C_{[19,6,9]_3}$. Extending $G_{[19,6,9]_3}$ to $G_{[20,7,9]_3}$ by adding ${\bf x}$, a ternary LCD $[20,7,9]$ code $C_{[20,7,9]_3}$ is constructed, where ${\bf x}=(1 1 0 2 0 0 1 1 0 0 0 0 0 1 1 0 2 2 2).$
By the MAGMA function BKLC, one can construct a ternary LCD $[20,12,6]$ code $C_{[20,12,6]_3}$.
\end{proof}






{\small
\begin{center}
 \begin{threeparttable}
\begin{tabular}{ccc}
\multicolumn{3}{c}{{\rm Table 6: The ternary LCD codes from Theorem \ref{LCD-(n-l,k-l)}}}\\
\hline
    \makebox[0.3\textwidth][c]{{\rm The\ linear\ code}\ $C_{[n,k,d]_3}^l$}& \makebox[0.3\textwidth][c]{The set $T$} &
    \makebox[0.3\textwidth][c]{The parameters for $C_T$} \\
    \hline\hline
$C^{0}_{[20,12,6]_3}$ &  $\backslash$& [20,12,6]\\

$C^{0}_{[21,12,6]_3} $& $\backslash$ & [21,12,6]\\

$C^{0}_{[21,15,4]_3}$ & $\backslash$ & [21,15,4]\\

$C^{0}_{[21,17,3]_3}$ & $\backslash$ & [21,17,3]\\

$C^{3}_{[25,9,11]_3} $& \{1,2,3\} & [22,6,11]\\

$C^{3}_{[25,13,8]_3}$ & \{1,2,8\} & [22,10,8]\\

$C^{0}_{[23,17,4]_3}$ &$ \backslash$& [23,17,4]\\

$C^{0}_{[24,18,4]_3} $& $\backslash$ & [24,18,4]\\

$C^{9}_{[44,29,8]_3} $& \{2,3,4,5,6,7,8,9,10\} & [35,20,8]\\

$C^{1}_{[40,29,6]_3}$ & \{2\} & [39,28,6]\\
\hline
\end{tabular}
\begin{tablenotes}
\footnotesize
\item [6] {\small We obtain some ternary LCD codes by Theorem \ref{LCD-(n-l,k-l)}, all the results are listed in Table 6, where the linear code $C_{[n,k,d]_2}^l$ in Table 6 is the best-known ternary $[n,k,d]$ code from MAGMA \cite{magma}.
It is worth noting that the code in MAGMA may not be in standard form, we need to use the MAGMA function StandardForm to act on it.}
\end{tablenotes}
\end{threeparttable}
\end{center}
}


\begin{cor}
$d_3^E(21,9)=8,d_3^E(22,10)=8,d_3^E(23,11)\in \{7,8\},d_3^E(24,12)\in \{7,8\}.$
\end{cor}

\begin{proof}
According to Table 6 and Table 7, there are ternary LCD $[21,9,8]$, $[22,10,8]$, $[23,11,7]$ and $[24,12,7]$ codes. From \cite{ter-11-19}, we know that $d_3^E(19,7)=d_3^E(20,8)=8$.
It follows from Corollary \ref{3-cor-leq} that $d_3^E(21,9)\leq \max\{d_3^E(20,8),d_3^E(19,7)\}=8$.
Similarly, $d_3^E(22,10)\leq 8$, $d_3^E(23,11)\leq 8$, $d_3^E(24,12)\leq 8$.
\end{proof}

\begin{remark}
Combining the Database \cite{codetables} with the above results, we give lower and upper bounds on minimal distance of ternary LCD codes with length $20\leq n\leq 25$, where the parameters for the ternary LCD of the length 20 can be found in \cite{ter-11-19}. All results are listed in Tables 8-9. In addition, the ternary LCD $[20,7,9]$ and $[20,12,6]$ codes have new parameters comparing to \cite{ter-11-19}.
\end{remark}


{\small
\begin{center}
 \begin{threeparttable}
\begin{tabular}{cccc}
\multicolumn{4}{c}{{\rm Table 7: The ternary LCD codes from Theorem \ref{Methods} and Theorem \ref{Methods-2}}}\\
\hline
    \makebox[0.15\textwidth][c]{{\rm The\ given}\ $C_{[n,k,d]_3}$}& \makebox[0.22\textwidth][c]{The vector ${\bf x}$ or ${\bf y}$} &
    \makebox[0.12\textwidth][c]{$C'$} &
    \makebox[0.12\textwidth][c]{References}\\
    \hline\hline
[19,6,9] &  (1102001100000110222) & [20,7,9] & {\rm Theorem}\ \ref{Methods} \\

[20,6,10] & (02112000000000121221) & [21,7,9] & {\rm Theorem}\ \ref{Methods} \\

[20,8,8] &  (12021210000020212222) & [21,9,8] & {\rm Theorem}\ \ref{Methods} \\

[20,5,11] &  (21112201000000010021) & [21,6,10] & {\rm Theorem}\ \ref{Methods} \\

[21,6,10] & (200002221000002020221) & [22,7,10] & {\rm Theorem}\ \ref{Methods} \\

[21,7,9] &  (222210000000002002222) & [22,8,9] & {\rm Theorem}\ \ref{Methods} \\

[21,12,6] &  (120000022222120122212) & [22,13,6] & {\rm Theorem}\ \ref{Methods} \\

[21,15,4] &  (010000111012101011211) & [22,16,4] & {\rm Theorem}\ \ref{Methods} \\

[21,17,3] &   (010021201100021201212) & [22,18,3] & {\rm Theorem}\ \ref{Methods} \\

[22,6,11] &   (1211021222210000112200) & [23,7,11] & {\rm Theorem}\ \ref{Methods} \\

[22,7,10] &   (1112202121010001202201) & [23,8,10] & {\rm Theorem}\ \ref{Methods} \\

[22,8,9] &   (2210022110000011120021) & [23,9,9] & {\rm Theorem}\ \ref{Methods} \\

[22,11,7] &   (0010021000012101211101) & [23,12,7] & {\rm Theorem}\ \ref{Methods} \\

[22,13,6] &    (2110210002112020111112) & [23,14,6] & {\rm Theorem}\ \ref{Methods} \\

[23,7,11] &    (10112100000000000102211) & [24,8,10] & {\rm Theorem}\ \ref{Methods} \\

[23,12,7] &     (00100211211100012121012) & [24,13,7] & {\rm Theorem}\ \ref{Methods} \\

[23,14,6] &     (22100010002211222202121) & [24,15,6] & {\rm Theorem}\ \ref{Methods} \\

[24,8,10] &     (020212110000000020010011) & [25,9,10] & {\rm Theorem}\ \ref{Methods} \\

[24,15,6] &    (102200010201120202000120) & [25,16,6] & {\rm Theorem}\ \ref{Methods} \\

[24,16,5] &    (010000010021012212210010) & [25,17,5] & {\rm Theorem}\ \ref{Methods} \\

[24,18,4] &    (121000002021110022212212) & [25,19,4] & {\rm Theorem}\ \ref{Methods} \\

[35,20,8] &       (12202110000000002102121001212020220) & [36,21,8] & {\rm Theorem}\ \ref{Methods} \\

[36,21,8] &        (100200212100000112202021222000012121) & [37,22,8] & {\rm Theorem}\ \ref{Methods} \\

[39,28,6] &         (021020210000200202112021120012101111112) & [40,29,6] & {\rm Theorem}\ \ref{Methods} \\

[22,10,8] &     (1101221102212211012200) & [22,11,7] & {\rm Theorem}\ \ref{Methods-2} \\

[22,13,6] &      (0020100000102202121012) & [22,14,5] & {\rm Theorem}\ \ref{Methods-2} \\

[23,14,6] &  (00020100000102202121012) & [23,15,5] & {\rm Theorem}\ \ref{Methods-2} \\

[24,15,6] &  (001121202101001122102021) & [24,16,5] & {\rm Theorem}\ \ref{Methods-2} \\

[25,9,10] &  (1112101001021000011020001) & [25,10,9] & {\rm Theorem}\ \ref{Methods-2}\\

[25,10,9] &  (1222121000002000212020021) & [25,11,8] & {\rm Theorem}\ \ref{Methods-2}\\
\hline
\end{tabular}
\begin{tablenotes}
\footnotesize
\item [7] {\small We obtain some new ternary LCD codes by construction methods in Theorems \ref{Methods} and \ref{Methods-2}. All results are listed in Table 7, where the given code $C_{[n,k,d]_3}$ in Table 7 is from Tables 6-7.}
\end{tablenotes}
\end{threeparttable}
\end{center}}











{\small
\begin{center}
 \begin{threeparttable}
\begin{tabular}{cccccccccccc}
\multicolumn{12}{c}{{\rm Table 8: Bound on the minimum distance of ternary LCD codes}}\\
\hline
    \makebox[0.05\textwidth][c]{$n\backslash k$}& \makebox[0.06\textwidth][c]{4} & \makebox[0.06\textwidth][c]{5}& \makebox[0.06\textwidth][c]{6}& \makebox[0.06\textwidth][c]{7}& \makebox[0.06\textwidth][c]{8}& \makebox[0.06\textwidth][c]{9}& \makebox[0.06\textwidth][c]{10}& \makebox[0.05\textwidth][c]{11}& \makebox[0.05\textwidth][c]{12}& \makebox[0.05\textwidth][c]{13}& \makebox[0.05\textwidth][c]{14}\\
    \hline\hline
20   & 12 & 11& 10& {\bf 9}& 8& 7-8& 7 &6 & {\bf6}& 5  & 4     \\

21   & 12 & 11-12& 10-11& 9-10& 8-9& 8& 7-8 &6-7 & 6& 5-6  & 4-5 \\

22   & 12-13 & 11-12& 11-12& 10-11& 9-10& 8-9& 8 & 7-8 & 6-7& 6  & 5-6   \\

23   & 13-14 &11-13& 11-12&  11-12& 10-11& 9-10& 8-9 &7-8 & 7-8& 6-7  & 6     \\

24   & 15 & 11-14& 11-13&  11-12& 10-11& 9-11& 8-10 &7-9 & 7-8& 7-8  & 6-7    \\

25   & 15-16 &  11-15& 11-14&  11-13&  10-12& 10-11& 9-11 &8-10 & 7-9& 7-8  & 6-8    \\
\hline
\end{tabular}
\begin{tablenotes}
\footnotesize
\item [8] The parameters in bold denote the corresponding code has
new parameters according to \cite{ter-11-19}.
\end{tablenotes}
\end{threeparttable}
\end{center}

}





{\small
\begin{center}
\begin{tabular}{cccccccccccc}
\multicolumn{12}{c}{{\rm Table 9: Bound on the minimum distance of ternary LCD codes}}\\
\hline
    \makebox[0.05\textwidth][c]{$n\backslash k$}& \makebox[0.06\textwidth][c]{15} & \makebox[0.06\textwidth][c]{16}& \makebox[0.06\textwidth][c]{17}& \makebox[0.06\textwidth][c]{18}& \makebox[0.06\textwidth][c]{19}& \makebox[0.06\textwidth][c]{20}& \makebox[0.06\textwidth][c]{21}& \makebox[0.05\textwidth][c]{22}& \makebox[0.05\textwidth][c]{23}& \makebox[0.05\textwidth][c]{24}& \makebox[0.05\textwidth][c]{25}\\
    \hline\hline
20   & 3-4& 3& 2& 2 & 2& 1 &  &  &  &   &       \\

21   & 4& 3& 3& 2 & 2& 1 & 1 &  &  &   &      \\

22   & 4-5& 4& 3& 3 & 2& 2& 2 & 1 &  &   &      \\

23   & 5-6& 4-5& 4& 3 & 3& 2& 2& 2 & 1 &   &        \\

24   & 6& 5-6& 4-5& 4 & 3& 3& 2& 2& 1 &  1 &    \\

25  & 6-7& 6 & 5-6& 4-5 & 4& 3& 3& 2& 2& 2  & 1 \\
\hline
\end{tabular}
\end{center}
}


\subsection{Lengths greater than 25}
In this subsection, we construct some new ternary LCD codes of lengths greater than 25 comparing with \cite{234-lcD}.

\begin{prop}
There are ternary LCD $[37,22,8]$ and $[40,29,6]$ codes.
\end{prop}

\begin{proof}
Applying Theorem \ref{Methods} to the generator matrix $G_{[35,20,8]_3}$ of the ternary LCD code $C_{[35,20,8]_3}$. Extending $G_{[35,20,8]_3}$ to $G_{[36,21,8]_3}$ by adding ${\bf x}$, a ternary LCD $[36,21,8]$ code $C_{[36,21,8]_3}$ is constructed, where ${\bf x}=(1 2 2 0 2 1 1 0 0 0 0 0 0 0 0 0 2 1 0 2 1 2 1 0 0 1 2 1 2 0 2 0 2 2 0)$.
Applying Theorem \ref{Methods} to the generator matrix $G_{[36,21,8]_3}$ of the ternary LCD code $C_{[36,21,8]_3}$. Extending $G_{[36,21,8]_3}$ to $G_{[37,22,8]_3}$ by adding ${\bf x}$, a ternary LCD $[37,22,8]$ code $C_{[37,22,8]_3}$ is constructed, where ${\bf x}=(1 0 0 2 0 0 2 1 2 1 0 0 0 0 0 1 1 2 2 0 2 0 2 1 2 2 2 0 0 0 0 1 2 1 2 1)$.
Applying Theorem \ref{Methods} to the generator matrix $G_{[39,28,6]_3}$ of the ternary LCD code $C_{[39,28,6]_3}$. Extending $G_{[39,28,6]_3}$ to $G_{[40,29,6]_3}$ by adding ${\bf x}$, a ternary LCD $[40,29,6]$ code is constructed, where ${\bf x}=(0 2 1 0 2 0 2 1 0 0 0 0 2 0 0 2 0 2 1 1 2 0 2 1 1 2 0 0 1 2 1 0 1 1 1 1 1 1 2)$.
\end{proof}

\begin{remark}
Although the method 1 in this paper is given by modifying the method in \cite{234-lcD}, we can construct some ternary LCD codes with better parameters compared with \cite{234-lcD}.
For example, the ternary LCD code of the length 37 with the dimension 22 has the minimum distance 8, while the ternary LCD code of the length 37 with the dimension 22 in \cite[Corollary 6.3]{234-lcD} has the minimum distance 7. That is to say, the ternary LCD code $C_{[37,22,8]_3}$ we obtained is also considered new.
The code $C_{[40,29,6]}$ also has better parameters compared with \cite[Corollary 6.3]{234-lcD}.
\end{remark}
















\section{New quaternary Hermitian LCD codes}
Let $\F_4=\{0,1,\omega,\omega^2\}$.
Let $d^H_4 (n, k)$ denote the largest minimum distance among all quaternary Hermitian LCD $[n, k]$ codes.
%The following lemma was proved in \cite[Theorem 1]{34-LCD}.
%\begin{lem}\label{lemma-leq-quaternary}
%Suppose that $2\leq k\leq n$. Then
%$d^H_4(n,k)\leq d^H_4(n,k-1).$
%\end{lem}


\begin{lem}\label{H-22,23,24}
There are quaternary Hermitian LCD $[22,12,7]$, $[23,13,7]$, $[24,14,7]$ and $[25,15,7]$ codes.
\end{lem}


\begin{proof}
By the MAGMA function BKLC, one can construct a quaternary $[25,15,7]$ code $C^{4}_{[25,15,7]_4}$ with 4-dimensional Hermitian hull.
The shortened code $(C^{4}_{[25,15,7]_4})_{T}$ on $T=\{1,2,3,4\}$ is a quaternary Hermitian LCD $[21,11,7]$ code $C_{[21,11,7]_4}$. Applying Theorem \ref{Methods} to the generator matrix $G_{[21,11,7]_4}$ of the ternary LCD code $C_{[21,11,7]_4}$. Extending $G_{[21,11,7]_4}$ to $G_{[22,12,7]_4}$ by adding ${\bf x}$, a quaternary Hermitian LCD $[22,12,7]$ code is constructed, where ${\bf x}=(1 \omega 1 1 0 0 0 0 0 \omega 0 \omega \omega^2 \omega \omega 0 0 \omega^2 \omega^2 \omega^2 0)$.
By the MAGMA function BKLC, one can construct a quaternary $[26,16,7]$ code $C^{3}_{[26,16,7]_4}$ with 3-dimensional Hermitian hull.
The shortened code $(C^{3}_{[26,16,7]_4})_{T}$ on $T=\{1,2,3\}$ is a quaternary Hermitian LCD $[23,13,7]$ code $C_{[23,13,7]_4}$.
Applying Theorem \ref{Methods} to the generator matrix $G_{[23,13,7]_4}$ of the quaternary Hermitian LCD code $C_{[23,13,7]_4}$. Extending $G_{[23,13,7]_4}$ to $G_{[24,14,7]_4}$ by adding ${\bf x}$, a quaternary Hermitian LCD $[24,14,7]$ code is constructed, where ${\bf x}=(0 \omega \omega^2 0 \omega 0 0 1 \omega 0 1 \omega \omega \omega^2 \omega^2 \omega^2 0 0 \omega 1 1\\ \omega \omega)$.
Applying Theorem \ref{Methods} to the generator matrix $G_{[24,14,7]_4}$ of the quaternary Hermitian LCD code $C_{[24,14,7]_4}$. Extending $G_{[24,14,7]_4}$ to $G_{[25,15,7]_4}$ by adding ${\bf x}$, a quaternary Hermitian LCD $[25,15,7]$ code is constructed, where ${\bf x}=(1 1 0 0 1 \omega^2 1 \omega \omega^2 0 \omega^2 \omega^2 1 0 0 \omega^2 \omega^2 \omega \omega \omega^2 0 1 \omega^2 \omega)$.
\end{proof}

\begin{remark}
Compared with Table 3 in \cite{H-25}, the quaternary Hermitian LCD codes in Lemma \ref{H-22,23,24} have better parameters than their parameters.
For example,
the quaternary Hermitian LCD code of the length 24 with the dimension 14 has the minimum distance 7, while the quaternary Hermitian LCD code of the length 24 with the dimension 14 in \cite[Table 3]{H-25} has the minimum distance 6. That is to say, the quaternary Hermitian LCD code $C_{[24,14,7]_4}$ we obtained is also considered new.
\end{remark}



We use $[[n,k,d;c]]$ to denote a quaternary entanglement-assisted quantum error correcting code (EAQECC) that encodes $k$ information qubits into $n$ channel qubits with the help of $c$ ebits, and $d$ is called the minimum distance of the EAQECC. EAQECCs were introduced by Brun et al. in \cite{EAQECC}, which include the standard quantum stabilizer codes as a special case.
It has shown that if there is a quaternary Hermitian
LCD $[n, k, d]$ code, then there is an EAQECC with parameters $[[n, k, d; n - k]]$ (see \cite{H-20}). Hence, as a consequence of Lemma \ref{H-22,23,24} and \cite[Lemma 3.3]{H-lCD-k=3}, we have that

\begin{cor}
For $(n,k,d)\in \{(22,12,7),(23,13,7),(24,14,7),(25,15,7)\}$ and a nonnegative integer $s$,
there is an EAQECC with parameters $[[n+\frac{4^k-1}{3}s,k,d+4^{k-1}s;n+\frac{4^k-1}{3}s-k]]$.
\end{cor}





\section{Conclusion}

In this paper, we introduced three methods for constructing LCD codes by modifying some known methods.
We used two counterexamples to show that the conjecture proposed by Bouyuklieva is invalid.
We constructed many new binary LCD codes, ternary LCD codes and quaternary Hermitian LCD codes, which improve the known lower bounds on the largest minimum weights. We also improved upper bounds on some ternary LCD codes.
It is worth mentioning that we constructed some LCD codes that cannot be constructed by the methods in \cite{234-lcD}.
Finally, as an application of quaternary Hermitian LCD codes, we found some EAQECCs with new parameters. We believe that our methods can produce more results for LCD codes.



\begin{thebibliography}{1}

\bibitem{AH-TLCD-1-10} M. Araya, M. Harada, On the classification of linear complementary dual codes, Discrete Math., 2019, \textbf {342}(1): 270-278.
\bibitem{H-lCD-k=3} M. Araya, M. Harada, On the classification of quaternary optimal Hermitian LCD codes, Cryptogr. Commun., 2022, https://doi.org/10.1007/s12095-021-00552-5.
\bibitem{AH-BLCD-17-24} M. Araya, M. Harada, On the minimum weights of binary linear complementary dual codes, Cryptogr. Commun., 2020, {\bf12}(2): 285-300.
\bibitem{AHS-BLCD} M. Araya, M. Harada, K. Saito, Characterization and classification of optimal LCD codes, Des. Codes Cryptogr., 2021, \textbf{89}: 617-640.
    \bibitem{ter-11-19} M. Araya, M. Harada, K. Saito, On the minimum weights of binary LCD codes and ternary LCD codes, 2020, https://arxiv.org/pdf/1908.08661.pdf.
\bibitem{H-LCD-IT} M. Araya, M. Harada, K. Saito, Quaternary Hermitian linear complementary dual codes, IEEE Trans. Information Theory, 2020, {\bf66}: 2751-2759.
        \bibitem{magma} W. Bosma, J. Cannon, C. Playoust, The Magma algebra system I: The user language, J. Symbolic Comput., 1997, {\bf24}: 235-265.
    \bibitem{B-LCD-40} S. Bouyuklieva, Optimal binary LCD codes, Des. Codes Cryptogr., 2021, \textbf{89}(11): 2445-2461.
\bibitem{EAQECC} T. Brun, I. Devetak, M. H. Hsieh, Correcting quantum errors with
entanglement, Science, 2006, {\bf314}(5789): 436-439.

\bibitem{lcd-appl} C. Carlet, S. Guilley, Complementary dual codes for counter-measures
to side-channel attacks, Adv. Math. Commun., 2016, \textbf{10}(1): 131-150.
%\bibitem{C-G-O-S-LCD-and-isometry} C. Carlet, C. G$\ddot{{\rm u}}$neri, F. $\ddot{{\rm O}}$zbudak, P. Sol\'e, A new concatenated type construction for LCD codes and isometry codes, Discrete Math., 2018, \textbf{341}(3): 830-835.
%\bibitem{self-orth-1}  A. R. Calderbank, E. M. Rains, P. M. Shor, N. J. A. Sloane, Quantum error correction via codes over GF(4), IEEE Trans. Information Theory, 1998, \textbf{44} (4): 1369-1387.
\bibitem{LCD-1} C. Carlet, S. Mesnager, C. Tang, Y. Qi, Euclidean and hermitian LCD MDS codes, Des. Codes Cryptogr., 2018, \textbf{86}(11): 2605-2618.
    \bibitem{LCD-new} C. Carlet, S. Mesnager, C. Tang, Y. Qi, New characterization and parametrization of LCD codes, IEEE Trans. Information Theory, 2019, {\bf65}(1): 39-49.
\bibitem{LCD-2} C. Carlet, S. Mesnager, C. Tang, Y. Qi, On $\sigma$-LCD codes, IEEE Trans. Information Theory, 2019, \textbf{65}(3): 1694-1704.
\bibitem{LCD-3} C. Carlet, S. Mesnager, C. Tang, Y. Qi, R. Pellikaan, Linear codes over $\F_q$ are equivalent to LCD codes for $q > 3$, IEEE Trans. Information Theory, 2018, \textbf {64}(4): 3010-3017.
%\bibitem{LCD-4} B. Chen, H. Liu, New constructions of MDS codes with complementary duals, IEEE Trans. Information Theory, 2018, \textbf{64}(8): 5776-5782.


%\bibitem{C-L-X-pi} H. Chen, S. Ling, C. Xing, Asymptotically good quantum codes exceeding the Ashikhmin-Litsyn-Tsfasman bound, IEEE Trans. Information Theory, 2001, \textbf{47}(5): 2055-2058.

\bibitem{LCD-5} S. T. Dougherty, J. Kim, B. $\ddot{{\rm O}}$zkaya, L. Sok, P. Sol\'e, The combinatorics of LCD codes: linear programming bound and orthogonal matrices, Int. J. Inf. Coding Theory, 2017, \textbf{4}(2-3): 116-128.
\bibitem{2-LCD-30} Q. Fu, R. Li, F. Fu, Y. Rao, On the construction of binary optimal LCD codes with short length, Int. J. Found. Comput. Sci., 2019, {\bf30}: 1237-1245.
\bibitem{bound-12} L. Galvez, J. L. Kim, N. Lee, Y. G. Roe, B. S. Won, Some bounds
on binary LCD codes, Cryptogr. Commun., 2018, \textbf{10}(4): 719-728.
\bibitem{codetables} M. Grassl, Bounds on the minimum distance of linear codes and quantum codes,
http://www.codetables.de. Accessed 4 March 2022.


\bibitem{LCD-6} C. G$\ddot{{\rm u}}$neri, B. $\ddot{{\rm O}}$zkaya, P. Sol\'e, Quasi-cyclic complementary dual codes, Finite Fields Appl., 2016, \textbf {42}: 67-80.
\bibitem{234-lcD} M. Harada, Construction of binary LCD codes, ternary LCD codes and quaternary Hermitian LCD codes. Des. Codes Cryptogr., 2021, {\bf89}(10): 2295-2312.
\bibitem{b-up-1} M. Harada, Existence of new extremal doubly-even codes and extremal singly--even codes. Des. Codes Cryptogr., 1996, {\bf8}(3): 273-283.
\bibitem{H-lcd-IJQI} M. Harada, Some optimal entanglement-assisted quantum codes constructed from quaternary Hermitian linear complementary dual codes. Int. J. Quantum Inf., 2019, {\bf17}: 1950053.
\bibitem{b-up-2} M. Harada, The existence of a self-dual $[70, 35, 12]$ code and formally self-dual codes. Finite Fields Appl., 1997, {\bf3}(2): 131-139.
\bibitem{b-up-3} M. Harada, H. Kimura, New extremal doubly-even $[64, 32, 12]$ codes. Des. Codes Cryptogr., 1995, {\bf6}(2): 91-96.
\bibitem{HS-BLCD-1-16} M. Harada, K. Saito, Binary linear complementary dual codes, Cryptogr. Commun., 2019, \textbf{11}(4): 677-696.
    \bibitem{34-LCD} M. Harada, K. Saito, Remark on subcodes of linear complementary dual codes, Inform. Process. Lett., 2020, {\bf159}(160): 105963.
\bibitem{Huffman} W. C. Huffman, V. Pless, Fundamentals of Error-Correcting Codes, Cambridge University Press, 2003.
    \bibitem{2-LCD-40} K. Ishizuka, K. Saito, Construction for both self-dual codes and LCD codes, Adv. Math. Commun., 2021, doi: 10.3934/amc.2021070.

\bibitem{b-up-4} J. L. Kim, New extremal self-dual codes of lengths 36, 38 and 58, IEEE Trans. Information Theory, 2001, {\bf47}(1): 386-393.
\bibitem{b-up-5} J. L. Kim, New self-dual codes over $GF(4)$ with the highest known minimum weights, IEEE Trans. Information Theory, 2001, {\bf47}(4): 1575-1580.
\bibitem{b-up-6} J. L. Kim, Y. Lee, Euclidean and Hermitian self-dual MDS codes over large finite fields, J. Comb. Theory, Ser. A, 2004, {\bf 105}(1): 79-95.

\bibitem{H-LCD-20} C. Lai, A. Ashikhmin, Linear programming bounds for entanglement-assisted quantum error-correcting codes by split weight enumerators, IEEE Trans. Information Theory, 2018, {\bf64}(1): 622-639.
\bibitem{b-up-7} H. Lee, Y. Lee, Construction of self-dual codes over finite rings $\Z_{p^m}$, J. Comb. Theory, Ser. A, 2008, {\bf 115}(3): 407-422.
    \bibitem{H-20} L. Lu, R. Li, L. Guo, Q. Fu, Maximal entanglement entanglement-assisted quantum codes constructed from linear codes. Quantum Information Processing, 2015, {\bf14}(1): 165-182.
    \bibitem{H-25} L. Lu, X. Zhan, S. Yang, H. Cao, Optimal quaternary Hermitian LCD codes, 2020, https://arxiv.org/pdf/2010.10166.pdf.
\bibitem{MacWilliams} F. J. MacWilliams, N. J. A. Sloane, The theory of Error Correcting Codes, Amsterdam. The Netherlands: North-Holland, 1977.
\bibitem{LCD-7} J. Massey, Linear codes with complementary duals, Discrete Math., 1992, \textbf{106-107}: 337-342.
%\bibitem{LCD-8} S. Mesnager, C. Tang, Y. Qi, Complementary dual algebraic geometry codes, IEEE Trans. Information Theory, 2018, \textbf{64}(4): 2390-2397.
%\bibitem{Sen-p-e-codes} Sendrier, N.: Finding the permutation between equivalent codes: the support splitting algorithm. IEEE Trans. Inf. Theory 46(4), 1193-1203 (2000).
%\bibitem{Sen-S-auto-group} N. Sendrier, G. Skersys, On the computation of the automorphism group of a linear code, in Proc. IEEE Int. Symp. Inf. Theory,
%Washington, DC, USA, Jun. (1)2001.
\bibitem{LCD-is-good} N. Sendrier, Linear codes with complementary duals meet the Gilbert-Varshamov bound, Discrete Math., 2004, \textbf{285}(1): 345-347.

\bibitem{LCD-9} M. Shi, D. Huang, L. Sok, P. Sol\'e, Double circulant LCD Codes over $\Z_4$, Finite Fields Appl., 2019, \textbf{58}: 133-144.
\bibitem{LCD-huang} M. Shi, D. Huang, L. Sok, P. Sol\'e, Double circulant self-dual and LCD codes over Galois rings, Adv. Math. Commun., 2019, {\bf13}(1): 171-183.
\bibitem{ACD-E-codes} M. Shi, S. Li, J. Kim, P. Sol\'e, LCD and ACD codes over a noncommutative non-unital ring with four elements, Cryptogr. Commun., 2022, {\bf14}(3): 627-640.
\bibitem{LCD-T-matric} M. Shi, F. $\ddot{{\rm O}}$zbudak, L. Xu, P. Sol\'e. LCD codes from tridiagonal Toeplitz matrices, Finite Fields Appl., 2021, \textbf{75}: 101892.
\bibitem{LCD-zhu} M. Shi, H. Zhu, L. Qian, L. Sok, P. Sol\'e, On self-dual and LCD double circulant and double negacirculant codes over $\F_q + u\F_q$, Cryptogr. Commun., 2020, {\bf12}(1): 53-70.
\bibitem{Lin-sok} L. Sok, On Hermitian LCD codes and their Gray image, Finite Fields Appl., 2020, \textbf{62}(2), 101623.
\bibitem{LCD-10} L. Sok, M. Shi, P. Sol\'e, Construction of optimal LCD codes over large finite fields, Finite Fields Appl., 2018, \textbf{50}: 138-153.
\end{thebibliography}




\end{document}

\begin{table}
\begin{center}
 Table 1: The binary LCD codes from Theorem \ref{LCD-(n-l,k-l)}
 \end{center}
{\small$$
\begin{array}{|c|c|c|}
\hline
{\rm The\ linear\ code}\ C_{[n,k,d]_2}^l & {\rm The\ set}\ T & {\rm The\ parameters\ for}\ C_T\\
\hline
C^{1}_{[30,11,10]_2} & \{1\} & [29,10,10]\\
\hline
C_{[36,21,7]_2}^6 & \{1,2,3,5,16,18\} & {\bf[30,15,7]}\\
\hline
C^{1}_{[32,21,6]_2} & \{1\} & [31,20,6]\\
\hline
C^{0}_{[33,22,6]_2} &\backslash & [33,22,6]\\
\hline
C^{1}_{[34,15,9]_2} & \{1\} & [33,14,9]\\
\hline
C^{0}_{[34,22,6]_2} & \backslash & [34,22,6]\\
\hline
C^{7}_{[43,27,7]_2} & \{1,2,3,5,9,17,25\} & {\bf[36,20,7]}\\
\hline
C^{3}_{[40,13,13]_2} & \{2,3,5\} & {\bf[37,10,13]}\\
\hline
C^{8}_{[45,29,7]_2} & \{1\} & {\bf[37,21,7]}\\
\hline
\end{array}
$$}
{\small We obtain some new binary LCD codes by Theorem \ref{LCD-(n-l,k-l)}, all the results are listed in Table 1, where the linear code $C_{[n,k,d]_2}^l$ in Table 1 is the best-known binary $[n,k,d]$ code from MAGMA \cite{magma}. It is worth noting that the code in MAGMA may not be in standard form, we need to use the MAGMA function StandardForm to act on it.}
\end{table}

\begin{table}
\begin{center}
Table 2: The binary LCD codes from Theorem \ref{Methods} and Theorem \ref{Methods-2}
\end{center}
{\small $$
\begin{array}{|c|c|c|c|}
\hline
{\rm The\ given}\ C_{[n,k,d]_2} & {\rm The\ vector}\ {\bf x}\ {\rm or}\ {\bf y}& C'& {\rm References}\\
\hline
[30,15,7] & (1 1 0 0 0 1 1 0 1 1 1 1 1 0 1 1 1 0 1 1 1 1 1 1 0 0 1 0 0 1) & {\bf[31,16,7]} & {\rm Theorem}\ \ref{Methods} \\
\hline
[31,16,7] & (1110101000110101100110111011001) & {\bf[32,17,7]}& {\rm Theorem}\ \ref{Methods} \\
\hline
[32,17,7] & (11010011101011010101111000110111) & {\bf[33,18,7]}& {\rm Theorem}\ \ref{Methods} \\
\hline
[33,14,9] & (111000000011111101010110010000110) & [34,15,9]& {\rm Theorem}\ \ref{Methods} \\
\hline
[34,18,8] & (1100000000000100001100101101010110) & {\bf[35,19,7]}&{\rm Theorem}\ \ref{Methods} \\
\hline
[34,22,6] &  (1111111111111111111111111111111111) & {\bf[35,23,6]}& {\rm Theorem}\ \ref{Methods} \\
\hline
[36,16,10] &  (110100000100000010011010100001101110) & [37,17,9]&{\rm Theorem}\ \ref{Methods} \\
\hline
[37,21,7] & (0011101010010010110110100111111001111) & {\bf[38,22,7]}&{\rm Theorem}\ \ref{Methods} \\
\hline
[38,10,14] &  (11010011110000000000000000000011101100) & [39,11,13]&{\rm Theorem}\ \ref{Methods} \\
\hline
[38,22,7] & (10110110101011110010110101001001010100) & {\bf[39,23,7]}&{\rm Theorem}\ \ref{Methods} \\
\hline
[29,10,10] & (00010101011110101001100010001) & [29,11,9] &{\rm Theorem}\ \ref{Methods-2} \\
\hline
[31,20,6] & (0000111101010000001110010011100) & [31,21,5] &{\rm Theorem}\ \ref{Methods-2} \\
\hline
\end{array}
$$}
We obtain some new binary LCD codes by construction methods in Theorems \ref{Methods} and \ref{Methods-2}, all the results are listed in Table 2, where the given code $C_{[n,k,d]_2}$ in Table 2 is from Tables 1-2.
\end{table}


\begin{table}
\begin{center}
 Table 3: The binary LCD codes from some inequalities
 \end{center}
{\small $$
\begin{array}{|c|c|c||c|c|c|}
\hline
{\rm The\ given}\ C &  C' & {\rm References}&{\rm The\ given}\ C &  C' & {\rm References}\\
\hline
[29,11,9] & [30,11,9] & {\rm Lemma\ \ref{lem-inequality-1}}&
[35,21,6] & [36,21,6] & {\rm Lemma\ \ref{lem-inequality-1}}\\
\hline
[29,11,9] & [31,11,10] & {\rm (3)\ of\ Lemma\ \ref{lemma-leq-2}}&[35,23,6] & [36,23,6] & {\rm Lemma\ \ref{lem-inequality-1}}\\
\hline
[31,16,7] & [31,15,7] & {\rm Lemma\ \ref{lemma-leq}}&[36,18,8]&[37,18,8]&{\rm Lemma\ \ref{lem-inequality-1}}\\
\hline
[31,16,7] & {\bf[32,16,8]} & {\rm (2)\ of\ Lemma\ \ref{lemma-leq-2}}&
[35,19,7]&{\bf[37,19,8]}&{\rm (3)\ of\ Lemma\ \ref{lemma-leq-2}}\\
\hline
[32,16,8] & {\bf[32,15,8]} & {\rm Lemma\ \ref{lemma-leq}}&
[36,20,7]&{\bf[37,20,8]}& {\rm (2)\ of\ Lemma\ \ref{lemma-leq-2}}\\
\hline
[31,21,5] & [32,21,5] & {\rm Lemma\ \ref{lem-inequality-1}}&
[36,23,6]&[37,23,6]&{\rm Lemma\ \ref{lem-inequality-1}}\\
\hline
[32,15,8] & [33,15,8] & {\rm Lemma\ \ref{lem-inequality-1}}&
[37,10,13]&{\bf[38,10,14]}&{\rm (2)\ of\ Lemma\ \ref{lemma-leq-2}}\\
\hline
[32,16,8] & {\bf[33,16,8]} & {\rm Lemma\ \ref{lem-inequality-1}}&
[37,17,9]&[38,17,9]&{\rm Lemma\ \ref{lem-inequality-1}}\\
\hline
[32,17,7] & {\bf[33,17,7]} & {\rm Lemma\ \ref{lem-inequality-1}}&
[37,19,8]&[38,19,8]&{\rm Lemma\ \ref{lem-inequality-1}}\\
\hline
[33,22,6] & {\bf[33,21,6]} & {\rm Lemma\ \ref{lemma-leq}}&[37,20,8]&{\bf[38,20,8]}&{\rm Lemma\ \ref{lem-inequality-1}}\\
\hline
[33,18,7] & {\bf[34,18,8]} & {\rm (2)\ of\ Lemma\ \ref{lemma-leq-2}}&
[37,21,7]&[38,21,7]&{\rm Lemma\ \ref{lem-inequality-1}}\\
\hline
[34,18,8] & {\bf[34,17,8]} & {\rm Lemma\ \ref{lemma-leq}}&[37,23,6]&[38,23,6]&{\rm Lemma\ \ref{lem-inequality-1}}\\
\hline
[34,17,8] & [34,16,8] & {\rm Lemma\ \ref{lemma-leq}}&
[38,10,14]&[39,10,14]&{\rm Lemma\ \ref{lem-inequality-1}}\\
\hline
[33,21,6] & [34,21,6] & {\rm Lemma\ \ref{lem-inequality-1}}&
[38,20,8]&[39,20,8]&{\rm Lemma\ \ref{lem-inequality-1}}\\
\hline
[34,17,8] & {\bf[35,17,8]} & {\rm Lemma\ \ref{lem-inequality-1}}&
[38,22,7]&{\bf[39,22,8]}&{\rm (2)\ of\ Lemma\ \ref{lemma-leq-2}}\\
\hline
[34,18,8] & {\bf[35,18,8]} & {\rm Lemma\ \ref{lem-inequality-1}}&[39,22,8]&{\bf[39,21,8]}&{\rm Lemma\ \ref{lemma-leq}}\\
\hline
[34,21,6] & {\bf[35,21,6]} & {\rm Lemma\ \ref{lem-inequality-1}}&
[39,10,14]&[40,10,14]&{\rm Lemma\ \ref{lem-inequality-1}}\\
\hline
[34,15,9] & [36,15,10] & {\rm (3)\ of\ Lemma\ \ref{lemma-leq-2}}&[39,11,13]&[40,11,13]&{\rm Lemma\ \ref{lem-inequality-1}}\\
\hline
[35,17,8] & [36,17,8] & {\rm Lemma\ \ref{lem-inequality-1}}&
[39,21,8]&{\bf[40,21,8]}&{\rm Lemma\ \ref{lem-inequality-1}}\\
\hline
[35,18,8] & {\bf[36,18,8]} & {\rm Lemma\ \ref{lem-inequality-1}}&[39,22,8]&{\bf[40,22,8]} &{\rm Lemma\ \ref{lem-inequality-1}}\\
\hline
[35,19,7] & [36,19,7] & {\rm Lemma\ \ref{lem-inequality-1}}&
[39,23,7]&[40,23,7]&{\rm Lemma\ \ref{lem-inequality-1}}\\
\hline
\end{array}
$$}
 we use some inequalities to obtain some new binary LCD codes, all the results are listed in Table 3, where the given code $C$ in Table 3 is from Tables 1-3.
\end{table}


\begin{table}
\begin{center}
Table 4: Bounds on the minimum diatance of binary LCD codes
\end{center}
{\small $$
\begin{array}{|cccccccccc|}
\hline
n\ k  & 9 & 10 &    11    & 12    & 13 & 14 & 15 & 16 & 17\\
\hline
29   & 10& 10 &{\bf9^\diamond} & 8& 8  & 8  & 6  &  6 & 6\\
\hline
30   & 11& 10 &{\bf9^\diamond}\!-\!10 & 9& 8  & 8  & {\bf7^*}  &  6 & 6\\
\hline
31   & 11& 10 &  {\bf10^\diamond} &10&9&8 & {\bf7^\diamond}\!-\!8& {\bf 7^*} & 6 \\
\hline
32   & 12& 11 &  10  & 10&9\!-\!10&8\!-\!9 & {\bf8^*}& {\bf8^*} & {\bf 7^*}\\
\hline
33  & 12& 12 &  {\bf11^*}  & 10&10& 9\!-\!10& {\bf8^\diamond}\!-\!9& {\bf 8^*} & {\bf 7^*}\!-\!8\\
\hline
34   & 13& 12 &  11\!-\!12  & 11&10& 10&  {\bf9^\diamond}\!-\!10& {\bf8^\diamond}\!-\!9 & {\bf8^*}\\
\hline
35   & 13\!-\!14& 12\!-\!13 &  12  & 12&10\!-\!11& 10& 9\!-\!10& 9\!-\!10 & {\bf8^*}\\
\hline
36   & 14& 12\!-\!14 &  12\!-\!13  & 12&10\!-\!12& 10\!-\!11& {\bf10^\diamond}& 10 & {\bf8^\diamond}\!-\!9\\
\hline
37   & 14& {\bf13^*}\!-\!14 &  12\!-\!14& 12\!-\!13&10\!-\!12& 10\!-\!12&10\!-\!11&10 & {\bf9^\diamond}\!-\!10\\
\hline
38  & 14\!-\!15& {\bf14^*} &  12\!-\!14& 12\!-\!14&11\!-\!12& 10\!-\!12&10\!-\!12&10\!-\!11 &{\bf9^\diamond}\!-\!10\\
\hline
39   & 14\!-\!16& {\bf14^\diamond}\!-\!15 &  {\bf13^\diamond}\!-\!14& 12\!-\!14&11\!-\!13& 11\!-\!12&10\!-\!12&10\!-\!12 & 10\!-\!11\\
\hline
40   & 15\!-\!16& {\bf14^\diamond}\!-\!16 &  {\bf13^\diamond}\!-\!15& 13\!-\!14&12\!-\!14& 11\!-\!13&10\!-\!12&10\!-\!12 & 10\!-\!12\\
\hline
\end{array}
$$}
\end{table}



\begin{table}
\begin{center}
Table 5: Bounds on the minimum diatance of binary LCD codes
\end{center}
{\footnotesize $$
\begin{array}{|cccccccccccccc|}
\hline
n\ k & 18 & 19 & 20 & 21 & 22 & 23 &    24    & 25    & 26 & 27 & 28 & 29 & 30\\
\hline
29   & 6  & 5& 4  & 4  & 4  & 3&     2    & 2  & 2  & 2  & 2  &  1 & \\
\hline
30   & 6  & 5\!-\!6& 5  & 4  & 4  & 4 &     3     & 2  & 2  & 2  & 2  &  1 & 1\\
\hline
31   & 6& 6& 6& {\bf5^\diamond} & 4  & 4  &  4    &   3   &2&2 & 2& 2 & 2 \\
\hline
32   & 6& 6& 6&  {\bf5^\diamond}\!-\!6& 5& 4 &  4     & 3\!-\!4&3&2 & 2& 2 & 2\\
\hline
33   & {\bf 7^*}& 6& 6&  {\bf6^*}&  6&  5 & 4  & 4&4& 3& 2& 2 & 2\\
\hline
34   & {\bf 8^*}& 6\!-\!7& 6& {\bf6^\diamond}&  6&  5\!-\!6 &  4  & 4&4& 3\!-\!4& 3&2 & 2\\
\hline
35   & {\bf 8^*}& {\bf 7^*}\!-\!8& 6\!-\!7& {\bf 6^*}&  6& {\bf 6^*} &   5  & 4&4& 4& 4& 3 & 2\\
\hline
36   &  {\bf8^*}& {\bf7^\diamond}\!-\!8& {\bf 7^*}\!-\!8&  {\bf6^\diamond}\!-\!7&  6&  {\bf6^\diamond} &  6  &  5&4& 4& 4& 3\!-\!4 & 3\\
\hline
37   &  {\bf8^\diamond}\!-\!9& {\bf 8^*}& {\bf 8^*}& {\bf 7^*}\!-\!8& 6\!-\!7&  {\bf6^\diamond} & 6&  5\!-\!6& 5& 4&4&4 & 4\\
\hline
38   & 9\!-\!10&  {\bf8^\diamond}\!-\!9&  {\bf8^*}&  {\bf7^\diamond}\!-\!8& {\bf7^*}\!-\!8& {\bf6^\diamond}\!-\!7 &   6& 6& 6&  5&4&4 & 4\\
\hline
39   & 10&  9\!-\!10& {\bf8^\diamond}\!-\!9& {\bf 8^*}& {\bf8^*}& {\bf7^*}\!-\!8 &   6\!-\!7&  6&6&  5\!-\!6& 5&4 & 4\\
\hline
40   & 10\!-\!11& 9\!-\!10&  9\!-\!10& {\bf 8^*}\!-\!9&{\bf8^*} & {\bf7^\diamond}\!-\!8 &  6\!-\!8&  6\!-\!7 & 6 &  6& 6&5 & 4\\
\hline
\end{array}
$$}
\end{table}

\begin{table}
\begin{center}
Table 6: The ternary LCD codes from Theorem \ref{LCD-(n-l,k-l)}
\end{center}
{\small $$
\begin{array}{|c|c|c|}
\hline
{\rm The\ linear\ code}\ C_{[n,k,d]_3}^l & {\rm The\ set}\ T & {\rm The\ parameters\ for}\ C_T\\
\hline
C^{0}_{[20,12,6]_3} &  \backslash& [20,12,6]\\
\hline
C^{0}_{[21,12,6]_3} & \backslash & [21,12,6]\\
\hline
C^{0}_{[21,15,4]_3} & \backslash & [21,15,4]\\
\hline
C^{0}_{[21,17,3]_3} & \backslash & [21,17,3]\\
\hline
C^{3}_{[25,9,11]_3} & \{1,2,3\} & [22,6,11]\\
\hline
C^{3}_{[25,13,8]_3} & \{1,2,8\} & [22,10,8]\\
\hline
C^{0}_{[23,17,4]_3} & \backslash& [23,17,4]\\
\hline
C^{0}_{[24,18,4]_3} & \backslash & [24,18,4]\\
\hline
C^{9}_{[44,29,8]_3} & \{2,3,4,5,6,7,8,9,10\} & [35,20,8]\\
\hline
C^{1}_{[40,29,6]_3} & \{2\} & [39,28,6]\\
\hline
\end{array}
$$}
\end{table}

\begin{table}
\begin{center}
Table 7: The ternary LCD codes from Theorem \ref{Methods} and Theorem \ref{Methods-2}
\end{center}
{\small $$
\begin{array}{|c|c|c|c|}
\hline
{\rm The\ given}\ C_{[n,k,d]_3} & {\rm The\ vector}\ {\bf x}\ {\rm or}\ {\bf y}& C'& {\rm References}\\
\hline
[19,6,9] &  (1102001100000110222) & [20,7,9] & {\rm Theorem}\ \ref{Methods} \\
\hline
[20,6,10] & (0 2 1 1 2 0 0 0 0 0 0 0 0 0 1 2 1 2 2 1 ) & [21,7,9] & {\rm Theorem}\ \ref{Methods} \\
\hline
[20,8,8] &  (12021210000020212222) & [21,9,8] & {\rm Theorem}\ \ref{Methods} \\
\hline
[20,5,11] &  (21112201000000010021) & [21,6,10] & {\rm Theorem}\ \ref{Methods} \\
\hline
[21,6,10] & (200002221000002020221) & [22,7,10] & {\rm Theorem}\ \ref{Methods} \\
\hline
[21,7,9] &  (222210000000002002222) & [22,8,9] & {\rm Theorem}\ \ref{Methods} \\
\hline
[21,12,6] &  (120000022222120122212) & [22,13,6] & {\rm Theorem}\ \ref{Methods} \\
\hline
[21,15,4] &  (010000111012101011211) & [22,16,4] & {\rm Theorem}\ \ref{Methods} \\
\hline
[21,17,3] &   (010021201100021201212) & [22,18,3] & {\rm Theorem}\ \ref{Methods} \\
\hline
[22,6,11] &   (1211021222210000112200) & [23,7,11] & {\rm Theorem}\ \ref{Methods} \\
\hline
[22,7,10] &   (1112202121010001202201) & [23,8,10] & {\rm Theorem}\ \ref{Methods} \\
\hline
[22,8,9] &   (2210022110000011120021) & [23,9,9] & {\rm Theorem}\ \ref{Methods} \\
\hline
[22,11,7] &   (0010021000012101211101) & [23,12,7] & {\rm Theorem}\ \ref{Methods} \\
\hline
[22,13,6] &    (2110210002112020111112) & [23,14,6] & {\rm Theorem}\ \ref{Methods} \\
\hline
[23,7,11] &    (10112100000000000102211) & [24,8,10] & {\rm Theorem}\ \ref{Methods} \\
\hline
[23,12,7] &     (00100211211100012121012) & [24,13,7] & {\rm Theorem}\ \ref{Methods} \\
\hline
[23,14,6] &     (22100010002211222202121) & [24,15,6] & {\rm Theorem}\ \ref{Methods} \\
\hline
[24,8,10] &     (020212110000000020010011) & [25,9,10] & {\rm Theorem}\ \ref{Methods} \\
\hline
[24,15,6] &    (102200010201120202000120) & [25,16,6] & {\rm Theorem}\ \ref{Methods} \\
\hline
[24,16,5] &    (010000010021012212210010) & [25,17,5] & {\rm Theorem}\ \ref{Methods} \\
\hline
[24,18,4] &    (121000002021110022212212) & [25,19,4] & {\rm Theorem}\ \ref{Methods} \\
\hline
[35,20,8] &       (12202110000000002102121001212020220) & [36,21,8] & {\rm Theorem}\ \ref{Methods} \\
\hline
[36,21,8] &        (100200212100000112202021222000012121) & [37,22,8] & {\rm Theorem}\ \ref{Methods} \\
\hline
[39,28,6] &         (021020210000200202112021120012101111112) & [40,29,6] & {\rm Theorem}\ \ref{Methods} \\
\hline
[22,10,8] &     (1101221102212211012200) & [22,11,7] & {\rm Theorem}\ \ref{Methods-2} \\
\hline
[22,13,6] &      (0020100000102202121012) & [22,14,5] & {\rm Theorem}\ \ref{Methods-2} \\
\hline
[23,14,6] &     (00020100000102202121012) & [23,15,5] & {\rm Theorem}\ \ref{Methods-2} \\
\hline
[24,15,6] &     (001121202101001122102021) & [24,16,5] & {\rm Theorem}\ \ref{Methods-2} \\
\hline
[25,9,10] &      (1112101001021000011020001) & [25,10,9] & {\rm Theorem}\ \ref{Methods-2} \\
\hline
[25,10,9] &       (1222121000002000212020021) & [25,11,8] & {\rm Theorem}\ \ref{Methods-2} \\
\hline
\end{array}
$$}
\end{table}





\begin{table}
\begin{center}
Table 8: Bound on the minimum distance of ternary LCD codes
\end{center}
\setlength{\arraycolsep}{3.2pt}
{\small $$
\begin{array}{|cccccccccccc|}
\hline
n\ k & 4 & 5 & 6 & 7 & 8 & 9 & 10 &    11    & 12    & 13 & 14   \\
\hline
20   & 12 & 11& 10& {\bf 9}& 8& 7\!-\!8& 7 &6 & {\bf6}& 5  & 4     \\
\hline
21   & 12 & 11\!-\!12& 10\!-\!11& 9\!-\!10& 8\!-\!9& 8& 7\!-\!8 &6\!-\!7 & 6& 5\!-\!6  & 4\!-\!5 \\
\hline
22   & 12\!-\!13 & 11\!-\!12& 11\!-\!12& 10\!-\!11& 9\!-\!10& 8\!-\!9& 8 & 7\!-\!8 & 6\!-\!7& 6  & 5\!-\!6   \\
\hline
23   & 13\!-\!14 &11\!-\!13& 11\!-\!12&  11\!-\!12& 10\!-\!11& 9\!-\!10& 8\!-\!9 &7\!-\!8 & 7\!-\!8& 6\!-\!7  & 6     \\
\hline
24   & 15 & 11\!-\!14& 11\!-\!13&  11\!-\!12& 10\!-\!11& 9\!-\!11& 8\!-\!10 &7\!-\!9 & 7\!-\!8& 7\!-\!8  & 6\!-\!7    \\
\hline
25   & 15\!-\!16 &  11\!-\!15& 11\!-\!14&  11\!-\!13&  10\!-\!12& 10\!-\!11& 9\!-\!11 &8\!-\!10 & 7\!-\!9& 7\!-\!8  & 6\!-\!8    \\
\hline
\end{array}
$$}
\end{table}


\begin{table}
\begin{center}
Table 9: Bound on the minimum distance of ternary LCD codes
\end{center}
{\small $$
\begin{array}{|cccccccccccc|}
\hline
n\ k & 15& 16& 17& 18 & 19 & 20 & 21 & 22 & 23 & 24 &    25     \\
\hline
20   & 3\!-\!4& 3& 2& 2 & 2& 1 &  &  &  &   &       \\
\hline
21   & 4& 3& 3& 2 & 2& 1 & 1 &  &  &   &      \\
\hline
22   & 4\!-\!5& 4& 3& 3 & 2& 2& 2 & 1 &  &   &      \\
\hline
23   & 5\!-\!6& 4\!-\!5& 4& 3 & 3& 2& 2& 2 & 1 &   &        \\
\hline
24   & 6& 5\!-\!6& 4\!-\!5& 4 & 3& 3& 2& 2& 1 &  1 &    \\
\hline
25  & 6\!-\!7& 6 & 5\!-\!6& 4\!-\!5 & 4& 3& 3& 2& 2& 2  & 1 \\
\hline

\end{array}
$$}
\end{table}



